\input{preamble}

% OK, start here.
%
\begin{document}

\title{Injectifs}


\maketitle

\phantomsection
\label{section-phantom}

\tableofcontents

\section{Introduction}
\label{section-introduction}

\noindent
Dans des chapitres ultérieurs, nous utiliserons l'existence d'injectifs et
de complexes K-injectifs pour étudier la cohomologie des faisceaux de modules
sur les sites annelés. Dans ce chapitre, nous expliquons comment construire
des injectifs et des complexes K-injectifs, d'abord pour les modules sur les
sites, puis, plus généralement, pour les catégories abéliennes de Grothendieck.

\medskip\noindent
Observons que nous savons déjà que la catégorie des groupes abéliens et la
catégorie des modules sur un anneau possèdent assez d'injectifs ; voir
Compléments d'algèbre, sections
\ref{more-algebra-section-abelian-groups} et
\ref{more-algebra-section-injectives-modules}





\section{L'argument de Baer pour les modules}
\label{section-baer}

\noindent
Il existe une autre manière, plus ensembliste, de montrer que tout $R$-module
$M$ se plonge dans un module injectif. Cette méthode construit le module
injectif comme une colimite transfinie de sommes amalgamées. Bien qu'elle soit
assez abstraite et plus compliquée que celle de
Compléments d'algèbre, section \ref{more-algebra-section-injectives-modules},
elle est aussi plus générale. Cette méthode semble remonter à Baer ; elle a
été reprise par Cartan et Eilenberg dans \cite{Cartan-Eilenberg}, puis par
Grothendieck dans \cite{Tohoku}. Grothendieck y montre que bien d'autres
catégories abéliennes possèdent assez d'injectifs. Nous reviendrons plus loin
au cas général (section \ref{section-grothendieck-categories}).

\medskip\noindent
Commençons par quelques remarques ensemblistes.
Soit $\{B_{\beta}\}_{\beta \in \alpha}$ un système inductif d'objets d'une
catégorie $\mathcal{C}$, indexé par un ordinal $\alpha$. Supposons que
$\colim_{\beta \in \alpha} B_\beta$ existe dans $\mathcal{C}$. Si $A$ est un
objet de $\mathcal{C}$, il existe alors une application naturelle
\begin{equation}
\label{equation-compare}
\colim_{\beta \in \alpha} \Mor_\mathcal{C}(A, B_\beta)
\longrightarrow
\Mor_\mathcal{C}(A, \colim_{\beta \in \alpha} B_\beta).
\end{equation}
En effet, un morphisme $A \to B_\beta$ donné pour un certain $\beta$ fournit
naturellement un morphisme de $A$ vers la colimite, par composition avec
$B_\beta \to \colim_{\beta \in \alpha} B_\alpha$.
Notons que la colimite de gauche est une colimite d'ensembles ! En général,
(\ref{equation-compare}) n'est ni injective ni surjective.

\begin{example}
\label{example-not-surjective}
Considérons la catégorie des ensembles. Posons $A = \mathbf{N}$ et prenons
pour $B_n = \{1, \ldots, n\}$ le système inductif indexé par les entiers
naturels, où $B_n \to B_m$, pour $n \leq m$, est l'application évidente.
Alors $\colim B_n = \mathbf{N}$ ; il existe donc une application
$A \to \colim B_n$ qui ne se factorise par aucune application $A \to B_m$,
pour tout $m$.
Par conséquent,
$\colim \Mor(A, B_n) \to \Mor(A, \colim B_n)$
n'est pas surjective.
\end{example}

\begin{example}
\label{example-not-injective}
Donnons ensuite un exemple où l'application n'est pas injective. Posons
$B_n = \mathbf{N}/\{1, 2, \ldots, n\}$, c'est-à-dire l'ensemble quotient de
$\mathbf{N}$ dans lequel les $n$ premiers éléments sont identifiés à un seul.
Il existe des applications naturelles $B_n \to B_m$ pour $n \leq m$ ; les
$\{B_n\}$ forment donc un système d'ensembles indexé par $\mathbf{N}$. On voit
aisément que $\colim B_n = \{*\}$ : c'est l'ensemble à un élément. Il s'ensuit
que $\Mor(A, \colim B_n)$ est un ensemble à un élément pour tout ensemble $A$.
Cependant, $\colim \Mor(A , B_n)$ n'est {\bf pas} un ensemble à un élément.
Considérons la famille d'applications $A \to B_n$ formée des projections
naturelles $\mathbf{N} \to \mathbf{N}/\{1, 2, \ldots, n\}$, et la famille
d'applications $A \to B_n$ qui envoient tout $A$ sur la classe de $1$.
Ces deux familles d'applications sont distinctes à chaque étape, donc distinctes
dans $\colim \Mor(A, B_n)$, mais elles induisent la même application
$A \to \colim B_n$.
\end{example}

\noindent
Néanmoins, si l'ensemble source est fini, (\ref{equation-compare}) est toujours
un isomorphisme.

\begin{lemma}
\label{lemma-out-of-finite}
Supposons que, dans (\ref{equation-compare}), $\mathcal{C}$ soit la catégorie
des ensembles et que $A$ soit un {\it ensemble fini}. Alors l'application est
une bijection.
\end{lemma}

\begin{proof}
Soit $f : A \to \colim B_\beta$. L'image de $f$ est finie ; notons ses éléments
$c_1, \ldots, c_r \in \colim B_\beta$. Par définition de la colimite, ils
proviennent tous d'éléments d'un même $B_\beta$, pour $\beta \in \alpha$
assez grand. On peut donc définir $\widetilde{f} : A \to B_\beta$ qui relève
$f$ à une étape finie. Cela prouve que (\ref{equation-compare}) est surjective.
Supposons maintenant que deux applications
$f : A \to B_\gamma, f' : A \to B_{\gamma'}$ définissent la même application
$A \to \colim B_\beta$. Chacun des éléments, en nombre fini, de $A$ a alors
la même image dans la colimite. Par définition de la colimite d'ensembles, il
existe $\beta \geq \gamma, \gamma'$ tel que les éléments de $A$ aient les mêmes images
dans $B_\beta$ par $f$ et $f'$. Cela prouve que (\ref{equation-compare}) est
injective.
\end{proof}

\noindent
Le cas le plus intéressant du lemme est celui où $\alpha = \omega$, autrement
dit où le système $\{B_\beta\}$ est un système
$\{B_n\}_{n \in \mathbf{N}}$ indexé par les entiers naturels, comme dans les
Exemples \ref{example-not-surjective} et \ref{example-not-injective}.
L'idée essentielle est que $A$ est « petit » relativement à la longue chaîne
de composées $B_1 \to B_2 \to \ldots$, de sorte que son morphisme doit se
factoriser par une étape finie. On trouvera une version plus générale de ce
lemme dans Ensembles, lemme \ref{sets-lemma-map-from-set-lifts}. Généralisons
maintenant ce résultat à la catégorie des modules.

\begin{definition}
\label{definition-small}
Soient $\mathcal{C}$ une catégorie, $I \subset \text{Arrows}(\mathcal{C})$
et $\alpha$ un ordinal. Un objet $A$ de $\mathcal{C}$ est dit
{\it $\alpha$-petit relativement à $I$} si, pour tout système
$\{B_\beta\}$ indexé par $\alpha$ dont les morphismes de transition
appartiennent à $I$, l'application (\ref{equation-compare}) est un
isomorphisme.
\end{definition}

\noindent
Dans la suite de cette section, nous nous restreindrons à la catégorie des
$R$-modules, où $R$ est un anneau commutatif fixé. Nous prendrons aussi pour
$I$ la famille des applications injectives, c'est-à-dire des
{\it monomorphismes} dans la catégorie des modules sur $R$. Dans ce cas, pour
tout système $\{B_\beta\}$ comme dans la définition, chacune des applications
$$
B_\beta \to \colim_{\beta \in \alpha} B_\beta
$$
est injective. Il s'ensuit que l'application (\ref{equation-compare}) est une
{\it injection}. On peut en fait considérer les $B_\beta$ comme des
sous-modules du module $B = \colim_{\beta \in \alpha} B_\beta$ ; on a alors
$B = \bigcup_{\beta \in \alpha} B_\beta$. Il n'y a pas abus de notation si
l'on identifie $B_\alpha$ à son image dans la colimite. Nous voulons maintenant
montrer que les modules sont toujours petits pour des ordinaux $\alpha$
« grands ».

\begin{proposition}
\label{proposition-modules-are-small}
Soient $R$ un anneau et $M$ un $R$-module. Soit $\kappa$ le cardinal de
l'ensemble des sous-modules de $M$. Si $\alpha$ est un ordinal dont la
cofinalité est strictement supérieure à $\kappa$, alors $M$ est
$\alpha$-petit relativement aux injections.
\end{proposition}

\begin{proof}
La démonstration est immédiate, mais considérons d'abord un cas particulier.
Si $M$ est fini, l'assertion est que, pour tout système inductif
$\{B_\beta\}$ à morphismes de transition injectifs et indexé par un ordinal
limite, toute application $M \to \colim B_\beta$ se factorise par l'un des
$B_\beta$. C'est précisément ce que nous avons démontré dans le
Lemme \ref{lemma-out-of-finite}.

\medskip\noindent
Passons au cas général. Il suffit de montrer que l'application
(\ref{equation-compare}) est surjective. Soit
$f : M \to \colim B_\beta$ une application. Considérons les sous-objets
$\{f^{-1}(B_\beta)\}$ de $M$, où $B_\beta$ est vu comme un sous-objet de la
colimite $B = \bigcup_\beta B_\beta$. Si l'un d'eux, disons
$f^{-1}(B_\beta)$, est égal à $M$, alors l'application se factorise par
$B_\beta$.

\medskip\noindent
Supposons donc, par l'absurde, que tous les $f^{-1}(B_\beta)$ soient des
sous-objets propres de $M$. Or nous savons que
$$
\bigcup f^{-1}(B_\beta) = f^{-1}\left(\bigcup B_\beta\right) = M.
$$
Par hypothèse, il y a au plus $\kappa$ sous-objets distincts de $M$ parmi les
$f^{-1}(B_\alpha)$. On peut donc trouver un sous-ensemble $S \subset \alpha$
de cardinal au plus $\kappa$ tel que, lorsque $\beta'$ parcourt $S$, les
$f^{-1}(B_{\beta'})$ parcourent \emph{tous} les $f^{-1}(B_\alpha)$.

\medskip\noindent
Cependant, $S$ possède un majorant $\widetilde{\alpha} < \alpha$, puisque
$\alpha$ est de cofinalité strictement supérieure à $\kappa$. En particulier,
tous les $f^{-1}(B_{\beta'})$, pour $\beta' \in S$, sont contenus dans
$f^{-1}(B_{\widetilde{\alpha}})$. Il en résulte que
$f^{-1}(B_{\widetilde{\alpha}}) = M$. Ainsi, l'application $f$ se factorise
par $B_{\widetilde{\alpha}}$.
\end{proof}

\noindent
Ce lemme nous permettra de déduire l'existence de nombreux injectifs.
Rappelons le critère de Baer.

\begin{lemma}[Critère de Baer]
\label{lemma-criterion-baer}
\begin{reference}
\cite[Theorem 1]{Baer}
\end{reference}
Soit $R$ un anneau. Un $R$-module $Q$ est injectif si et seulement si, dans
tout diagramme commutatif
$$
\xymatrix{
\mathfrak{a} \ar[d] \ar[r] &  Q \\
R \ar@{-->}[ru]
}
$$
où $\mathfrak{a} \subset R$ est un idéal, la flèche pointillée existe.
\end{lemma}

\begin{proof}
Il s'agit de l'équivalence entre (1) et (3) dans
Compléments d'algèbre, lemme
\ref{more-algebra-lemma-characterize-injective-bis};
notons que la démonstration qui y est donnée est élémentaire (elle n'utilise ni
les groupes $\text{Ext}$, ni l'existence d'injectifs ou de projectifs dans la
catégorie des $R$-modules).
\end{proof}

\noindent
Si $M$ est un $R$-module, nous pouvons en général disposer d'un diagramme
partiellement complété comme dans le Lemme \ref{lemma-criterion-baer}. Nous
pouvons y former la \emph{somme amalgamée}
$$
\xymatrix{
\mathfrak{a} \ar[d]  \ar[r] &  M \ar[d] \\
R \ar[r] &  R \oplus_{\mathfrak{a}} M.
}
$$
Ici, le morphisme vertical est injectif et le diagramme est commutatif. Le point
essentiel est que nous pouvons prolonger $\mathfrak{a} \to M$ à $R$ \emph{si}
nous remplaçons $M$ par le module plus grand
$R \oplus_{\mathfrak{a}} M$.

\medskip\noindent
L'idée fondamentale de l'argument de Baer consiste à répéter cette procédure
un nombre transfini de fois. Pour cela, étant donné un $R$-module $M$, nous
commençons par définir l'immense somme amalgamée suivante
\begin{equation}
\label{equation-huge-diagram}
\vcenter{
\xymatrix{
\bigoplus_{\mathfrak a}
\bigoplus_{\varphi \in \Hom_R(\mathfrak a, M)}
\mathfrak{a} \ar[r] \ar[d] & M \ar[d] \\
\bigoplus_{\mathfrak a}
\bigoplus_{\varphi \in \Hom_R(\mathfrak a, M)}
R \ar[r] &  \mathbf{M}(M).
}
}
\end{equation}
La flèche horizontale supérieure envoie l'élément $a \in \mathfrak a$ du
facteur correspondant à $\varphi$ sur l'élément $\varphi(a) \in M$. La flèche
verticale de gauche envoie simplement $a \in \mathfrak a$, dans le facteur
correspondant à $\varphi$, sur l'élément $a \in R$ du facteur correspondant à
$\varphi$. Les propriétés fondamentales de cette construction sont énoncées
dans le lemme suivant.

\begin{lemma}
\label{lemma-construction}
Soit $R$ un anneau.
\begin{enumerate}
\item La construction $M \mapsto (M \to \mathbf{M}(M))$ est fonctorielle
en $M$.
\item L'application $M \to \mathbf{M}(M)$ est injective.
\item Pour tout idéal $\mathfrak{a}$ et tout morphisme de $R$-modules
$\varphi : \mathfrak a \to M$, il existe un morphisme de $R$-modules
$\varphi' : R \to \mathbf{M}(M)$ tel que le diagramme
$$
\xymatrix{
\mathfrak{a} \ar[d] \ar[r]_\varphi &  M \ar[d] \\
R \ar[r]^{\varphi'} & \mathbf{M}(M)
}
$$
soit commutatif.
\end{enumerate}
\end{lemma}

\begin{proof}
Les assertions (2) et (3) résultent immédiatement de la construction. Pour
voir (1), soit $\chi : M \to N$ un morphisme de $R$-modules. Nous affirmons
qu'il existe un diagramme commutatif canonique
$$
\xymatrix{
\bigoplus_{\mathfrak a}
\bigoplus_{\varphi \in \Hom_R(\mathfrak a, M)}
\mathfrak{a} \ar[r] \ar[d] \ar[rrd] & M \ar[rrd]^\chi \\
\bigoplus_{\mathfrak a}
\bigoplus_{\varphi \in \Hom_R(\mathfrak a, M)}
R \ar[rrd] & &
\bigoplus_{\mathfrak a}
\bigoplus_{\psi \in \Hom_R(\mathfrak a, N)}
\mathfrak{a} \ar[r] \ar[d] & N \\
& & \bigoplus_{\mathfrak a}
\bigoplus_{\psi \in \Hom_R(\mathfrak a, N)}
R
}
$$
qui induit l'application voulue $\mathbf{M}(M) \to \mathbf{M}(N)$. La flèche
centrale dirigée vers l'est-sud-est envoie le facteur $\mathfrak a$
correspondant à $\varphi$, par $\text{id}_{\mathfrak a}$, sur le facteur
$\mathfrak a$ correspondant à
$\psi = \chi \circ \varphi$. Il en va de même de la flèche inférieure dirigée
vers l'est-sud-est. Nous omettons les détails.
\end{proof}

\noindent
Nous allons maintenant appliquer le foncteur $\mathbf{M}$ un nombre transfini
de fois. Pour chaque ordinal $\alpha$, nous définissons un foncteur
$\mathbf{M}_\alpha$ sur la catégorie des $R$-modules, ainsi qu'une injection
naturelle $N \to \mathbf{M}_\alpha(N)$, par récurrence transfinie. D'abord,
$\mathbf{M}_1 = \mathbf{M}$ est le foncteur défini ci-dessus. Soit ensuite un
ordinal $\alpha$, et supposons $\mathbf{M}_{\alpha'}$ défini pour
$\alpha' < \alpha$. Si $\alpha$ admet un prédécesseur immédiat
$\widetilde{\alpha}$, posons
$$
\mathbf{M}_\alpha = \mathbf{M} \circ \mathbf{M}_{\widetilde{\alpha}}.
$$
Sinon, c'est-à-dire si $\alpha$ est un ordinal limite, posons
$$
\mathbf{M}_{\alpha}(N) =
\colim_{\alpha' < \alpha} \mathbf{M}_{\alpha'}(N).
$$
Il est clair, par exemple par récurrence, que les
$\mathbf{M}_{\alpha}(N)$ forment un système inductif indexé par les ordinaux ;
cette définition a donc bien un sens.

\begin{theorem}
\label{theorem-baer-grothendieck}
Soit $\kappa$ le cardinal de l'ensemble des idéaux de $R$, et soit $\alpha$
un ordinal de cofinalité strictement supérieure à $\kappa$. Alors
$\mathbf{M}_\alpha(N)$ est un $R$-module injectif, et
$N \to \mathbf{M}_\alpha(N)$ est un plongement fonctoriel dans un injectif.
\end{theorem}

\begin{proof}
D'après le critère de Baer, Lemme \ref{lemma-criterion-baer}, il suffit de
montrer que, pour tout idéal $\mathfrak{a} \subset R$, toute application
$f : \mathfrak{a} \to \mathbf{M}_\alpha(N)$ se prolonge en une application
$R \to \mathbf{M}_\alpha(N)$. Or, puisque $\alpha$ est un ordinal limite, nous
savons que
$$
\mathbf{M}_{\alpha}(N) =
\colim_{\beta < \alpha} \mathbf{M}_{\beta}(N),
$$
donc, par la Proposition \ref{proposition-modules-are-small}, nous obtenons
$$
\Hom_R(\mathfrak{a}, \mathbf{M}_{\alpha}(N)) =
\colim_{\beta < \alpha} \Hom_R(\mathfrak a, \mathbf{M}_{\beta}(N)).
$$
Cela signifie en particulier qu'il existe $\beta' < \alpha$ tel que $f$ se
factorise par le sous-module $\mathbf{M}_{\beta'}(N)$, sous la forme
$$
f : \mathfrak{a} \to \mathbf{M}_{\beta'}(N) \to
\mathbf{M}_{\alpha}(N).
$$
Cependant, d'après la propriété fondamentale du foncteur $\mathbf{M}$, voir le
Lemme \ref{lemma-construction}, assertion (3), l'application
$\mathfrak{a} \to \mathbf{M}_{\beta'}(N)$ se prolonge en
$$
R \to \mathbf{M}( \mathbf{M}_{\beta'}(N)) =
\mathbf{M}_{\beta' + 1}(N),
$$
et ce dernier objet se plonge dans $\mathbf{M}_{\alpha}(N)$, car
$\beta' + 1 < \alpha$ puisque $\alpha$ est un ordinal limite. Par conséquent,
$f$ se prolonge à $\mathbf{M}_{\alpha}(N)$.
\end{proof}




\section{G-modules}
\label{section-G-modules}

\noindent
Nous verrons plus loin (Algèbre différentielle graduée, section
\ref{dga-section-modules-noncommutative}) que la catégorie des modules sur une
algèbre admet des plongements fonctoriels dans des injectifs. La construction est
exactement la même que celle de
Compléments d'algèbre, section
\ref{more-algebra-section-injectives-modules}.

\begin{lemma}
\label{lemma-G-modules}
Soient $G$ un groupe topologique et $R$ un anneau. La catégorie
$\text{Mod}_{R, G}$ des $R\text{-}G$-modules, voir Cohomologie étale,
définition \ref{etale-cohomology-definition-G-module-continuous}, admet des
plongements fonctoriels dans des injectifs. Il en est ainsi, en particulier, de la
catégorie des $G$-modules discrets.
\end{lemma}

\begin{proof}
La remarque qui précède le lemme montre que la catégorie
$\text{Mod}_{R[G]}$ admet des plongements fonctoriels dans des injectifs. Considérons le
foncteur d'oubli
$v : \text{Mod}_{R, G} \to \text{Mod}_{R[G]}$.
Ce foncteur est pleinement fidèle, transforme les applications injectives en
applications injectives et possède un adjoint à droite, à savoir
$$
u : M \mapsto u(M) = \{x \in M \mid \text{le stabilisateur de }x\text{ est ouvert}\}
$$
Comme $v(M) = 0 \Rightarrow M = 0$, on conclut par
Homologie, lemme \ref{homology-lemma-adjoint-functorial-injectives}.
\end{proof}



\section{Faisceaux abéliens sur un espace}
\label{section-abelian-sheaves-space}


\begin{lemma}
\label{lemma-abelian-sheaves-space}
Soit $X$ un espace topologique. La catégorie des faisceaux abéliens sur $X$
possède assez d'injectifs. Elle admet même des plongements fonctoriels dans des
injectifs.
\end{lemma}

\begin{proof}
Pour tout groupe abélien $A$, notons $j : A \to J(A)$ le plongement fonctoriel
dans un injectif construit dans
Compléments d'algèbre, section
\ref{more-algebra-section-injectives-modules}.
Soit $\mathcal{F}$ un faisceau abélien sur $X$. D'après Faisceaux, exemple
\ref{sheaves-example-sheaf-product-pointwise}, la donnée
$$
\mathcal{I} : U \mapsto
\mathcal{I}(U) = \prod\nolimits_{x\in U} J(\mathcal{F}_x)
$$
est un faisceau abélien. Il existe une application canonique
$\mathcal{F} \to \mathcal{I}$ qui envoie $s \in \mathcal{F}(U)$ sur
$\prod_{x \in U} j(s_x)$, où $s_x \in \mathcal{F}_x$ désigne le germe de $s$
en $x$. Cette application est injective ; voir, par exemple, Faisceaux, lemme
\ref{sheaves-lemma-sheaf-subset-stalks}.

\medskip\noindent
Il reste à démontrer l'assertion suivante : si une donnée $x \mapsto I_x$
associe à chaque point $x \in X$ un groupe abélien injectif, alors le faisceau
$\mathcal{I} : U \mapsto \prod_{x \in U} I_x$ est injectif. Notons que
$$
\mathcal{I} = \prod\nolimits_{x \in X} i_{x, *}I_x
$$
est le produit des faisceaux gratte-ciel $i_{x, *}I_x$ (pour la notation, voir
Faisceaux, section \ref{sheaves-section-skyscraper-sheaves}). On a
$$
\Mor_{\textit{Ab}}(\mathcal{F}_x, I_x)
=
\Mor_{\textit{Ab}(X)}(\mathcal{F}, i_{x, *}I_x).
$$
voir Faisceaux, lemme \ref{sheaves-lemma-stalk-skyscraper-adjoint}. Il est donc
clair que chaque $i_{x, *}I_x$ est injectif. L'injectivité de $\mathcal{I}$
résulte alors de
Homologie, lemme \ref{homology-lemma-product-injectives}.
\end{proof}









\section{Faisceaux de modules sur un espace annelé}
\label{section-sheaves-modules-space}


\begin{lemma}
\label{lemma-sheaves-modules-space}
Soit $(X, \mathcal{O}_X)$ un espace annelé ; voir Faisceaux, section
\ref{sheaves-section-ringed-spaces}. La catégorie des faisceaux de
$\mathcal{O}_X$-modules sur $X$ possède assez d'injectifs. Elle admet même des
plongements fonctoriels dans des injectifs.
\end{lemma}

\begin{proof}
Pour tout anneau $R$ et tout $R$-module $M$, notons
$j : M \to J_R(M)$ le plongement fonctoriel dans un injectif construit dans
Compléments d'algèbre, section
\ref{more-algebra-section-injectives-modules}.
Soit $\mathcal{F}$ un faisceau de $\mathcal{O}_X$-modules sur $X$. D'après
Faisceaux, exemples \ref{sheaves-example-sheaf-product-pointwise} et
\ref{sheaves-example-sheaf-product-pointwise-algebraic-structure}, la donnée
$$
\mathcal{I} : U \mapsto
\mathcal{I}(U) = \prod\nolimits_{x\in U} J_{\mathcal{O}_{X, x}}(\mathcal{F}_x)
$$
est un faisceau abélien. Il existe une application canonique
$\mathcal{F} \to \mathcal{I}$ qui envoie $s \in \mathcal{F}(U)$ sur
$\prod_{x \in U} j(s_x)$, où $s_x \in \mathcal{F}_x$ désigne le germe de $s$
en $x$. Cette application est injective ; voir, par exemple, Faisceaux, lemme
\ref{sheaves-lemma-sheaf-subset-stalks}.

\medskip\noindent
Il reste à démontrer l'assertion suivante : si une donnée $x \mapsto I_x$
associe à chaque point $x \in X$ un $\mathcal{O}_{X, x}$-module injectif, alors
le faisceau $\mathcal{I} : U \mapsto \prod_{x \in U} I_x$ est injectif.
Notons que
$$
\mathcal{I} = \prod\nolimits_{x \in X} i_{x, *}I_x
$$
est le produit des faisceaux gratte-ciel $i_{x, *}I_x$ (pour la notation, voir
Faisceaux, section \ref{sheaves-section-skyscraper-sheaves}). On a
$$
\Hom_{\mathcal{O}_{X, x}}(\mathcal{F}_x, I_x)
=
\Hom_{\mathcal{O}_X}(\mathcal{F}, i_{x, *}I_x).
$$
voir Faisceaux, lemme \ref{sheaves-lemma-stalk-skyscraper-adjoint}. Il est donc
clair que chaque $i_{x, *}I_x$ est un $\mathcal{O}_X$-module injectif (voir
Homologie, lemme \ref{homology-lemma-adjoint-preserve-injectives}, ou raisonner
directement). L'injectivité de $\mathcal{I}$ résulte alors de
Homologie, lemme \ref{homology-lemma-product-injectives}.
\end{proof}













\section{Préfaisceaux abéliens sur une catégorie}
\label{section-injectives-presheaves}

\noindent
Soit $\mathcal{C}$ une catégorie. Rappelons que cela signifie que
$\Ob(\mathcal{C})$ est un ensemble. Considérons, d'une part, les préfaisceaux
abéliens sur $\mathcal{C}$, voir Sites, section
\ref{sites-section-presheaves}, et, d'autre part, les familles de groupes
abéliens indexées par les éléments de $\Ob(\mathcal{C})$ ; autrement dit, les
préfaisceaux sur la catégorie discrète dont l'ensemble sous-jacent d'objets est
$\Ob(\mathcal{C})$. Notons simplement $\Ob(\mathcal{C})$ cette catégorie
discrète. Il existe un foncteur naturel
$$
i : \Ob(\mathcal{C}) \longrightarrow \mathcal{C}
$$
et, par conséquent, un foncteur naturel de restriction, ou foncteur d'oubli,
$$
v = i^p :
\textit{PAb}(\mathcal{C})
\longrightarrow
\textit{PAb}(\Ob(\mathcal{C}))
$$
à comparer avec Sites, section \ref{sites-section-functoriality-PSh}. Nous
noterons les préfaisceaux sur $\mathcal{C}$ par $B$ et les préfaisceaux sur
$\Ob(\mathcal{C})$ par $A$.

\medskip\noindent
Il existe aussi deux foncteurs, à savoir $i_p$ et ${}_pi$, qui associent un
préfaisceau abélien sur $\mathcal{C}$ à un préfaisceau abélien sur
$\Ob(\mathcal{C})$ ; voir Sites, sections
\ref{sites-section-functoriality-PSh} et
\ref{sites-section-more-functoriality-PSh}. Nous utiliserons ici
$u = {}_pi$, qui, dans le cas présent, est défini par
$$
uA(U) = \prod\nolimits_{U' \to U} A(U').
$$
Ainsi, un élément est une famille $(a_\phi)_\phi$, où $\phi$ parcourt tous les
morphismes de $\mathcal{C}$ de but $U$. L'application de restriction de $uA$
correspondant à $g : V \to U$ envoie l'élément $(a_\phi)_\phi$ sur
$(a_{g \circ \psi})_\psi$.

\medskip\noindent
Il existe une application surjective canonique $vuA \to A$ et une application
injective canonique $B \to uvB$. Nous laissons au lecteur le soin de montrer
que
$$
\Mor_{\textit{PAb}(\mathcal{C})}(B, uA)
=
\Mor_{\textit{PAb}(\Ob(\mathcal{C}))}(vB, A).
$$
dans ce cas simple ; le cas général est traité dans Sites, section
\ref{sites-section-functoriality-PSh}. Ainsi, le couple $(u, v)$ est un exemple
de couple de foncteurs adjoints ; voir Catégories, section
\ref{categories-section-adjoint}.

\medskip\noindent
Nous pouvons maintenant énumérer les faits suivants concernant la situation
ci-dessus.
\begin{enumerate}
\item Les foncteurs $u$ et $v$ sont exacts. Cela résulte de leur description
explicite donnée ci-dessus.
\item En particulier, le foncteur $v$ transforme les applications injectives
en applications injectives.
\item La catégorie $\textit{PAb}(\Ob(\mathcal{C}))$ possède assez d'injectifs.
\item Il existe même un plongement fonctoriel dans un injectif
$A \mapsto \big(A \to J(A)\big)$ comme dans
Homologie, définition \ref{homology-definition-functorial-injective-embedding}.
On peut en effet prendre pour $J(A)$ le préfaisceau
$U\mapsto J(A(U))$, où $J(-)$ est le foncteur construit dans
Compléments d'algèbre, section
\ref{more-algebra-section-injectives-modules}
pour l'anneau $\mathbf{Z}$.
\end{enumerate}
En réunissant ces faits, on obtient la procédure suivante pour plonger les
objets $B$ de $\textit{PAb}(\mathcal{C})$ dans un objet injectif :
$B \to uJ(vB)$. Voir
Homologie, lemme \ref{homology-lemma-adjoint-functorial-injectives}.

\begin{proposition}
\label{proposition-presheaves-injectives}
Les préfaisceaux abéliens sur une catégorie admettent un plongement fonctoriel
dans un injectif.
\end{proposition}

\begin{proof}
Voir la discussion ci-dessus.
\end{proof}

\section{Faisceaux abéliens sur un site}
\label{section-injectives-sheaves}

\noindent
Soit $\mathcal{C}$ un site. Dans cette section, nous démontrons que la
catégorie des faisceaux abéliens sur $\mathcal{C}$ a assez d'injectifs.

\medskip\noindent
Notons
$i : \textit{Ab}(\mathcal{C}) \longrightarrow \textit{PAb}(\mathcal{C})$
le foncteur d'oubli des faisceaux abéliens vers les préfaisceaux abéliens.
Notons
${}^\# : \textit{PAb}(\mathcal{C}) \longrightarrow \textit{Ab}(\mathcal{C})$
le foncteur de faisceautisation. Rappelons que ${}^\#$ est adjoint à gauche
de $i$, que ${}^\#$ est exact et que $i\mathcal{F}^\# = \mathcal{F}$
pour tout faisceau abélien $\mathcal{F}$. Enfin, notons
$\mathcal{G} \to J(\mathcal{G})$ le plongement canonique dans un
préfaisceau injectif construit à la
section \ref{section-injectives-presheaves}.

\medskip\noindent
Pour tout faisceau $\mathcal{F}$ de $\textit{Ab}(\mathcal{C})$ et tout
ordinal $\beta$, nous définissons un faisceau $J_\beta(\mathcal{F})$ par
récurrence transfinie. Posons $J_0(\mathcal{F}) = \mathcal{F}$ et
$J_1(\mathcal{F}) = J(i\mathcal{F})^\#$. La faisceautisation de
l'application canonique $i\mathcal{F} \to J(i\mathcal{F})$ donne une
application fonctorielle
$$
\mathcal{F} \longrightarrow J_1(\mathcal{F})
$$
qui est injective puisque $\#$ est exact. Posons
$J_{\alpha + 1}(\mathcal{F}) = J_1(J_\alpha(\mathcal{F}))$. On obtient ainsi
des applications injectives canoniques
$J_\alpha(\mathcal{F}) \to J_{\alpha + 1}(\mathcal{F})$. Pour un ordinal
limite $\beta$, nous posons
$$
J_\beta(\mathcal{F}) = \colim_{\alpha < \beta} J_\alpha(\mathcal{F}).
$$
Remarquons qu'il s'agit d'une colimite filtrante. Ainsi, pour tous ordinaux
$\alpha < \beta$, on dispose d'une application injective
$J_\alpha(\mathcal{F}) \to J_\beta(\mathcal{F})$.

\begin{lemma}
\label{lemma-map-into-next-one}
Avec les notations ci-dessus, supposons que
$\mathcal{G}_1 \to \mathcal{G}_2$ soit une application injective de faisceaux
abéliens sur $\mathcal{C}$. Soit $\alpha$ un ordinal et soit
$\mathcal{G}_1 \to J_\alpha(\mathcal{F})$ un morphisme de faisceaux. Il existe
un morphisme $\mathcal{G}_2 \to J_{\alpha + 1}(\mathcal{F})$ tel que le
diagramme suivant soit commutatif
$$
\xymatrix{
\mathcal{G}_1 \ar[d] \ar[r] & \mathcal{G}_2 \ar[d] \\
J_{\alpha}(\mathcal{F}) \ar[r] & J_{\alpha + 1}(\mathcal{F}) }
$$
\end{lemma}

\begin{proof}
En effet, l'application $i\mathcal{G}_1 \to i\mathcal{G}_2$ est injective ;
par conséquent, $i\mathcal{G}_1 \to iJ_\alpha(\mathcal{F})$ se prolonge en
$i\mathcal{G}_2 \to J(iJ_\alpha(\mathcal{F}))$, ce qui donne l'application
voulue après application du foncteur de faisceautisation.
\end{proof}

\noindent
Ce lemme affirme qu'en un certain sens le système
$\{J_{\alpha}(\mathcal{F})\}$ est un plongement injectif de $\mathcal{F}$.
Bien entendu, on ne peut pas prendre la limite sur tous les $\alpha$, car
ceux-ci forment une classe et non un ensemble. L'idée est cependant qu'il
n'est pas nécessaire de vérifier l'injectivité sur toutes les injections
$\mathcal{G}_1 \to \mathcal{G}_2$, compte tenu du lemme suivant.

\begin{lemma}
\label{lemma-map-into-smaller}
Supposons que les $\mathcal{G}_i$, $i\in I$, forment un ensemble de faisceaux
abéliens sur $\mathcal{C}$. Il existe un ordinal $\beta$ tel que, pour tout
faisceau $\mathcal{F}$, tout $i\in I$ et toute application
$\varphi : \mathcal{G}_i \to J_\beta(\mathcal{F})$, il existe un
$\alpha < \beta$ tel que $ \varphi $ se factorise par
$J_\alpha(\mathcal{F})$.
\end{lemma}

\begin{proof}
On se ramène au cas d'un seul faisceau $\mathcal{G}$ en prenant la somme
directe de tous les $\mathcal{G}_i$.

\medskip\noindent
Considérons les ensembles
$$
S = \coprod\nolimits_{U \in \Ob(\mathcal{C})} \mathcal{G}(U).
$$
et
$$
T_\beta
=
\coprod\nolimits_{U \in \Ob(\mathcal{C})} J_\beta(\mathcal{F})(U)
$$
Les applications de transition entre les ensembles $T_\beta$ sont injectives.
Si la cofinalité de $\beta$ est assez grande, alors
$T_\beta = \colim_{\alpha < \beta} T_\alpha$ ; voir
Sites, lemme \ref{sites-lemma-colimit-over-ordinal-sections}. Un morphisme
$\mathcal{G} \to J_\beta(\mathcal{F})$ se factorise par
$J_\alpha(\mathcal{F})$ si et seulement si l'application associée
$S \to T_\beta$ se factorise par $T_\alpha$. D'après
Ensembles, lemme \ref{sets-lemma-map-from-set-lifts}, si la cofinalité de $\beta$
est strictement supérieure au cardinal de $S$, la conclusion du lemme est
vraie. Le lemme résulte donc du fait qu'il existe des ordinaux de cofinalité
arbitrairement grande ; voir
Ensembles, proposition \ref{sets-proposition-exist-ordinals-large-cofinality}.
\end{proof}

\noindent
Rappelons que, pour un objet $X$ de $\mathcal{C}$, on note $\mathbf{Z}_X$ le
préfaisceau de groupes abéliens $\Gamma(U, \mathbf{Z}_X) =
\oplus_{U \to X} \mathbf{Z}$ ; voir
Modules sur les sites, section \ref{sites-modules-section-free-abelian-presheaf}.
Le faisceau associé à ce préfaisceau est noté $\mathbf{Z}_X^\#$ ; voir
Modules sur les sites, section \ref{sites-modules-section-free-abelian-sheaf}.
Il se caractérise par la propriété
\begin{equation}
\label{equation-free-sheaf-on}
\Mor_{\textit{Ab}(\mathcal{C})}(\mathbf{Z}_X^\#, \mathcal{G})
=
\mathcal{G}(X)
\end{equation}
où l'élément $\varphi$ du membre de gauche est envoyé sur
$\varphi(1 \cdot \text{id}_X)$ dans celui de droite. Ces faisceaux permettent
de caractériser les faisceaux abéliens injectifs.

\begin{lemma}
\label{lemma-characterize-injectives}
Soit $\mathcal{J}$ un faisceau de groupes abéliens ayant la propriété
suivante : pour tout $X\in \Ob(\mathcal{C})$, tout sous-faisceau abélien
$\mathcal{S} \subset \mathbf{Z}_X^\#$ et tout morphisme
$\varphi : \mathcal{S} \to \mathcal{J}$, il existe un morphisme
$\mathbf{Z}_X^\# \to \mathcal{J}$ prolongeant $\varphi$. Alors $\mathcal{J}$
est un faisceau injectif de groupes abéliens.
\end{lemma}

\begin{proof}
Soit $\mathcal{F} \to \mathcal{G}$ une application injective de faisceaux
abéliens et soit $\varphi : \mathcal{F} \to \mathcal{J}$ un morphisme.
En raisonnant comme dans la preuve de
Compléments d'algèbre, lemme \ref{more-algebra-lemma-injective-abelian}, on voit qu'il
suffit de démontrer que, si $\mathcal{F} \not = \mathcal{G}$, on peut trouver
un faisceau abélien $\mathcal{F}'$ tel que
$\mathcal{F} \subset \mathcal{F}' \subset \mathcal{G}$, l'inclusion
$\mathcal{F} \subset \mathcal{F}'$ étant stricte, et tel que $\varphi$ se
prolonge à $\mathcal{F}'$. Pour trouver $\mathcal{F}'$, choisissons un objet
$X$ de $\mathcal{C}$ pour lequel l'inclusion
$\mathcal{F}(X) \subset \mathcal{G}(X)$ est stricte, puis
$s \in \mathcal{G}(X)$, $s \not \in \mathcal{F}(X)$. Soit
$\psi : \mathbf{Z}_X^\# \to \mathcal{G}$ le morphisme correspondant à la
section $s$ par (\ref{equation-free-sheaf-on}). Posons
$\mathcal{S} = \psi^{-1}(\mathcal{F})$. Par hypothèse, le morphisme
$$
\mathcal{S} \xrightarrow{\psi} \mathcal{F} \xrightarrow{\varphi} \mathcal{J}
$$
se prolonge en un morphisme $\varphi' : \mathbf{Z}_X^\# \to \mathcal{J}$.
Remarquons que $\varphi'$ annule le noyau de $\psi$ (puisque c'est le cas de
$\varphi$). Ainsi, $\varphi'$ induit un morphisme
$\varphi'' : \Im(\psi) \to \mathcal{J}$ qui coïncide par construction avec
$\varphi$ sur l'intersection $\mathcal{F} \cap \Im(\psi)$. Les morphismes
$\varphi$ et $\varphi''$ se recollent donc en un prolongement de $\varphi$ au
sous-faisceau strictement plus grand
$\mathcal{F}' = \mathcal{F} + \Im(\psi)$.
\end{proof}

\begin{theorem}
\label{theorem-sheaves-injectives}
La catégorie des faisceaux de groupes abéliens sur un site a assez
d'injectifs. Plus précisément, il existe un plongement injectif fonctoriel ;
voir Homology, définition
\ref{homology-definition-functorial-injective-embedding}.
\end{theorem}

\begin{proof}
Soit $\mathcal{G}_i$, $i \in I$, un ensemble de faisceaux abéliens tel que
tout sous-faisceau de tout $\mathbf{Z}_X^\#$ soit l'un des $\mathcal{G}_i$.
Appliquons le lemme \ref{lemma-map-into-smaller} à cette famille et obtenons
ainsi un ordinal $\beta$. Nous affirmons que, pour tout faisceau de groupes
abéliens $\mathcal{F}$, l'application
$\mathcal{F} \to J_\beta(\mathcal{F})$ plonge $\mathcal{F}$ dans un injectif.
Remarquons que, par construction, l'association
$\mathcal{F} \mapsto \big(\mathcal{F} \to J_\beta(\mathcal{F})\big)$ est bien
fonctorielle.

\medskip\noindent
D'après le lemme \ref{lemma-characterize-injectives}, il suffit, pour prouver
cette affirmation, de prolonger tout morphisme
$\gamma : \mathcal{G} \to J_\beta(\mathcal{F})$ défini sur un sous-faisceau
$\mathcal{G}$ d'un $\mathbf{Z}_X^\#$ à tout $\mathbf{Z}_X^\#$. Or, d'après le
lemme \ref{lemma-map-into-smaller}, l'application $\gamma$ se relève dans
$J_\alpha(\mathcal{F})$ pour un certain $\alpha < \beta$. Enfin, le
lemme \ref{lemma-map-into-next-one} fournit le prolongement voulu de $\gamma$
en un morphisme à valeurs dans
$J_{\alpha + 1}(\mathcal{F}) \to J_\beta(\mathcal{F})$.
\end{proof}






\section{Modules sur un site annelé}
\label{section-sheaves-modules}

\noindent
Soit $\mathcal{C}$ un site et soit $\mathcal{O}$ un faisceau d'anneaux sur
$\mathcal{C}$. Par analogie avec
Compléments d'algèbre, section \ref{more-algebra-section-injectives-modules},
cherchons à démontrer qu'il existe assez de $\mathcal{O}$-modules injectifs.
Choisissons tout d'abord un plongement injectif
$$
\bigoplus\nolimits_{U, \mathcal{I}}
j_{U!}\mathcal{O}_U/\mathcal{I}
\longrightarrow
\mathcal{J}
$$
où $\mathcal{J}$ est un faisceau abélien injectif, dont l'existence résulte de
la section précédente. La somme directe porte ici sur tous les objets $U$ de
$\mathcal{C}$ et tous les sous-$\mathcal{O}$-modules
$\mathcal{I} \subset j_{U!}\mathcal{O}_U$. Pour les foncteurs de restriction
et de prolongement par $0$ associés au foncteur de localisation
$j_U : \mathcal{C}/U \to \mathcal{C}$, voir
Modules sur les sites, section \ref{sites-modules-section-localize}.

\medskip\noindent
Pour tout faisceau de $\mathcal{O}$-modules $\mathcal{F}$, notons
$$
\mathcal{F}^\vee
=
\SheafHom(\mathcal{F}, \mathcal{J})
$$
muni de sa structure naturelle de $\mathcal{O}$-module.
Insérer ici une future référence au Hom interne.
Nous aurons également besoin d'une résolution plate canonique d'un faisceau de
$\mathcal{O}$-modules. Nous la construisons comme suit : pour tout
$\mathcal{O}$-module $\mathcal{F}$, notons
$$
F(\mathcal{F})
=
\bigoplus\nolimits_{U \in \Ob(\mathcal{C}), s \in \mathcal{F}(U)}
j_{U!}\mathcal{O}_U.
$$
C'est un faisceau plat de $\mathcal{O}$-modules, muni d'une surjection canonique
$F(\mathcal{F}) \to \mathcal{F}$ ; voir
Modules sur les sites, lemme \ref{sites-modules-lemma-module-quotient-flat}.
En outre, la construction $\mathcal{F} \mapsto F(\mathcal{F})$ est
fonctorielle en $\mathcal{F}$.

\begin{lemma}
\label{lemma-vee-exact-sheaves}
Le foncteur $\mathcal{F} \mapsto \mathcal{F}^\vee$ est exact.
\end{lemma}

\begin{proof}
Cela résulte de ce que $\mathcal{J}$ est un faisceau abélien injectif.
\end{proof}

\noindent
Il existe une application canonique d'évaluation
$ev : \mathcal{F} \to (\mathcal{F}^\vee)^\vee$ : étant donné
$x \in \mathcal{F}(U)$, on définit
$ev(x) \in (\mathcal{F}^\vee)^\vee =
\SheafHom(\mathcal{F}^\vee, \mathcal{J})$ comme l'application
$\varphi \mapsto \varphi(x)$.

\begin{lemma}
\label{lemma-ev-injective-sheaves}
Pour tout $\mathcal{O}$-module $\mathcal{F}$, l'application d'évaluation
$ev : \mathcal{F} \to (\mathcal{F}^\vee)^\vee$ est injective.
\end{lemma}

\begin{proof}
On le vérifie à l'aide de la définition de $\mathcal{J}$. En effet, si
$s \in \mathcal{F}(U)$ est non nul, soit
$j_{U!}\mathcal{O}_U \to \mathcal{F}$ l'application de
$\mathcal{O}$-modules qui lui correspond par adjonction, et soit
$\mathcal{I}$ son noyau. Il existe une application non nulle
$\mathcal{F} \supset j_{U!}\mathcal{O}_U/\mathcal{I} \to \mathcal{J}$ qui
n'annule pas $s$. Comme $\mathcal{J}$ est un faisceau abélien injectif, cette
application sous-jacente de $\mathcal{O}$-modules se prolonge en une application
$\varphi : \mathcal{F} \to \mathcal{J}$.
Alors $ev(s)(\varphi) = \varphi(s) \not = 0$, ce qu'il fallait démontrer.
\end{proof}

\noindent
La surjection canonique $F(\mathcal{F}) \to \mathcal{F}$ de
$\mathcal{O}$-modules donne, d'après ce qui précède, une injection canonique de
$\mathcal{O}$-modules
$$
(\mathcal{F}^\vee)^\vee \longrightarrow (F(\mathcal{F}^\vee))^\vee.
$$
Posons $J(\mathcal{F}) = (F(\mathcal{F}^\vee))^\vee$. La composée de $ev$
avec l'application affichée donne
$\mathcal{F} \to J(\mathcal{F})$, fonctoriellement en $\mathcal{F}$.

\begin{lemma}
\label{lemma-JM-injective-sheaves}
Soit $\mathcal{O}$ un faisceau d'anneaux. Pour tout $\mathcal{O}$-module
$\mathcal{F}$, le $\mathcal{O}$-module $J(\mathcal{F})$ est injectif.
\end{lemma}

\begin{proof}
Il faut montrer que le foncteur
$\Hom_\mathcal{O}(\mathcal{G}, J(\mathcal{F}))$ est exact. Remarquons que
\begin{eqnarray*}
\Hom_\mathcal{O}(\mathcal{G}, J(\mathcal{F}))
& = &
\Hom_\mathcal{O}(\mathcal{G}, (F(\mathcal{F}^\vee))^\vee) \\
& = &
\Hom_\mathcal{O}
(\mathcal{G}, \SheafHom(F(\mathcal{F}^\vee), \mathcal{J})) \\
& = &
\Hom(\mathcal{G} \otimes_\mathcal{O} F(\mathcal{F}^\vee), \mathcal{J})
\end{eqnarray*}
La conclusion résulte donc de ce que $F(\mathcal{F}^\vee)$ est plat et
$\mathcal{J}$ injectif.
\end{proof}

\begin{theorem}
\label{theorem-sheaves-modules-injectives}
Soit $\mathcal{C}$ un site et soit $\mathcal{O}$ un faisceau d'anneaux sur
$\mathcal{C}$. La catégorie des faisceaux de $\mathcal{O}$-modules sur ce site
a assez d'injectifs. Plus précisément, il existe un plongement injectif
fonctoriel ; voir Homology, définition
\ref{homology-definition-functorial-injective-embedding}.
\end{theorem}

\begin{proof}
Cela résulte de la discussion de cette section.
\end{proof}

\begin{proposition}
\label{proposition-presheaves-modules}
Soit $\mathcal{C}$ une catégorie et soit $\mathcal{O}$ un préfaisceau
d'anneaux sur $\mathcal{C}$. La catégorie
$\textit{PMod}(\mathcal{O})$ des préfaisceaux de $\mathcal{O}$-modules possède
des plongements injectifs fonctoriels.
\end{proposition}

\begin{proof}
On pourrait procéder comme dans la discussion de la
section \ref{section-injectives-presheaves}. Nous utiliserons plutôt le
théorème précédent. Munissons $\mathcal{C}$ de la structure de site pour
laquelle les recouvrements d'un objet $U$ sont tous les singletons
$\{f : V \to U\}$ où $f$ est un isomorphisme. Nous omettons la vérification
qu'il s'agit bien d'un site. Pour cette topologie, un faisceau est la même chose
qu'un préfaisceau (preuve omise). Le théorème s'applique donc.
\end{proof}








\section{Plongement des catégories abéliennes}
\label{section-embedding}

\noindent
Dans cette section, nous montrons qu'une catégorie abélienne se plonge dans la
catégorie des faisceaux abéliens sur un site ayant assez de points. Le site
sera celui décrit dans le lemme suivant.

\begin{lemma}
\label{lemma-site-abelian-category}
Soit $\mathcal{A}$ une catégorie abélienne. Posons
$$
\text{Cov} = \{\{f : V \to U\} \mid f\text{ est surjectif}\}.
$$
Alors $(\mathcal{A}, \text{Cov})$ est un site ; voir
Sites, définition \ref{sites-definition-site}.
\end{lemma}

\begin{proof}
Selon nos conventions sur les catégories, $\Ob(\mathcal{A})$ est un ensemble.
Un isomorphisme est un morphisme surjectif, la composée de morphismes
surjectifs est surjective, et le changement de base d'un morphisme surjectif
dans $\mathcal{A}$ est surjectif ; voir
Homology, lemme \ref{homology-lemma-epimorphism-universal-abelian-category}.
\end{proof}

\noindent
Soit $\mathcal{A}$ une catégorie préadditive. Dans ce cas, le plongement de
Yoneda $\mathcal{A} \to \textit{PSh}(\mathcal{A})$, $X \mapsto h_X$, se
factorise par un foncteur $\mathcal{A} \to \textit{PAb}(\mathcal{A})$.

\begin{lemma}
\label{lemma-embedding}
Soit $\mathcal{A}$ une catégorie abélienne et soit
$\mathcal{C} = (\mathcal{A}, \text{Cov})$ le site défini au
lemme \ref{lemma-site-abelian-category}. Alors $X \mapsto h_X$ définit un
foncteur exact et pleinement fidèle
$$
\mathcal{A} \longrightarrow \textit{Ab}(\mathcal{C}).
$$
De plus, le site $\mathcal{C}$ a assez de points.
\end{lemma}

\begin{proof}
Supposons que $f : V \to U$ soit un morphisme surjectif de $\mathcal{A}$ et
posons $K = \Ker(f)$. Rappelons que
$V \times_U V = \Ker((f, -f) : V \oplus V \to U)$ ; voir
Homology, exemple \ref{homology-example-fibre-product-pushouts}. Il existe en
particulier une injection $K \oplus K \to V \times_U V$. Soient
$p, q : V \times_U V \to V$ les deux projections. Le morphisme
$p - q : V \times_U V \to V$ vérifie $f \circ (p  - q) = 0$ ; ainsi,
$p - q$ se factorise par $K \to V$. Notons
$c : V \times_U V \to K$ le morphisme ainsi
obtenu. Comme la composée $K \oplus K \to V \times_U V \to K$ est
surjective, $c$ est surjectif. Par conséquent,
$$
V \times_U V \xrightarrow{p - q} V \to U \to 0
$$
est une suite exacte de $\mathcal{A}$. Ainsi, pour un objet $X$ de
$\mathcal{A}$, la suite
$$
0 \to
\Hom_\mathcal{A}(U, X) \to
\Hom_\mathcal{A}(V, X) \to
\Hom_\mathcal{A}(V \times_U V, X)
$$
est une suite exacte de groupes abéliens ; voir
Homology, lemme \ref{homology-lemma-check-exactness}. Cela signifie que $h_X$
satisfait à la condition de faisceau sur $\mathcal{C}$.

\medskip\noindent
Le foncteur est pleinement fidèle d'après
Catégories, lemme \ref{categories-lemma-yoneda}. Il est exact à gauche entre
catégories abéliennes d'après
Homology, lemme \ref{homology-lemma-check-exactness}. Pour montrer qu'il est
exact à droite, soit $X \to Y$ un morphisme surjectif de $\mathcal{A}$. Soit
$U$ un objet de $\mathcal{A}$ et soit
$s \in h_Y(U) = \Mor_\mathcal{A}(U, Y)$ une section de $h_Y$ au-dessus de
$U$. D'après
Homology, lemme \ref{homology-lemma-epimorphism-universal-abelian-category},
la projection $U \times_Y X \to U$ est surjective. Ainsi,
$\{V = U \times_Y X \to U\}$ est un recouvrement de $U$ tel que $s|_V$ se
relève en une section de $h_X$. Ceci montre que $h_X \to h_Y$ est une
surjection de faisceaux abéliens ; voir
Sites, lemme \ref{sites-lemma-mono-epi-sheaves}.

\medskip\noindent
Le site $\mathcal{C}$ a assez de points d'après
Sites, proposition \ref{sites-proposition-criterion-points}.
\end{proof}

\begin{remark}
\label{remark-embedding}
Le théorème de plongement de Freyd--Mitchell affirme qu'il existe, de toute
catégorie abélienne $\mathcal{A}$ vers la catégorie des modules sur un anneau,
un foncteur exact et pleinement fidèle. Le lemme \ref{lemma-embedding} n'est
pas tout à fait aussi fort. Son résultat convient toutefois au Champs Project,
puisque nous devons de toute façon étudier en détail les faisceaux de groupes
abéliens sur des sites. En outre, la « chasse au diagramme » fonctionne dans la
catégorie des faisceaux abéliens sur $\mathcal{C}$, par exemple en travaillant
avec des sections au-dessus des objets, ou au niveau des germes en utilisant le
fait que $\mathcal{C}$ a assez de points. Pour déduire le théorème de plongement
de Freyd--Mitchell du lemme \ref{lemma-embedding}, voir la
remarque \ref{remark-embedding-freyd}.
\end{remark}

\begin{remark}
\label{remark-embedding-big}
Si $\mathcal{A}$ est une « grande » catégorie abélienne, c'est-à-dire si
$\mathcal{A}$ possède une classe d'objets, le lemme \ref{lemma-embedding} ne
s'applique pas. Dans ce cas, pour tout ensemble d'objets
$E \subset \Ob(\mathcal{A})$, il existe une sous-catégorie abélienne pleine
$\mathcal{A}' \subset \mathcal{A}$ telle que $\Ob(\mathcal{A}')$ soit un
ensemble et $E \subset \Ob(\mathcal{A}')$. On peut alors appliquer le
lemme \ref{lemma-embedding} à $\mathcal{A}'$. Cela permet de prouver que les
résultats fondés sur une chasse au diagramme restent valables dans
$\mathcal{A}$.
\end{remark}

\begin{remark}
\label{remark-embedding-freyd}
Soit $\mathcal{C}$ un site. La catégorie $\textit{Ab}(\mathcal{C})$ a assez
d'injectifs ; voir le théorème \ref{theorem-sheaves-injectives}. Lorsque
$\mathcal{C}$ a assez de points, cela est immédiat, car $p_*I$ est un faisceau
injectif si $I$ est un $\mathbf{Z}$-module injectif et $p$ un point. De plus,
$\textit{Ab}(\mathcal{C})$ possède un cogénérateur (détails omis). Le
lemme \ref{lemma-embedding} donne donc un plongement exact et pleinement fidèle
$\mathcal{A} \to \mathcal{B}$, où $\mathcal{B}$ possède un cogénérateur et
assez d'injectifs. En l'appliquant à $\mathcal{A}^{opp}$, on obtient un foncteur
exact et pleinement fidèle
$i : \mathcal{A} \to \mathcal{D} = \mathcal{B}^{opp}$, où $\mathcal{D}$ a
assez de projectifs et possède un générateur. Ainsi, $\mathcal{D}$ possède un
générateur projectif $P$. Posons $R = \Mor_\mathcal{D}(P, P)$. Alors
$$
\mathcal{A} \longrightarrow \text{Mod}_R, \quad
X \longmapsto \Hom_\mathcal{D}(P, X).
$$
On vérifie que ce foncteur est exact et pleinement fidèle. Autrement dit, on
retrouve le théorème de Freyd--Mitchell mentionné dans la
remarque \ref{remark-embedding}.
\end{remark}

\begin{remark}
\label{remark-embed-exact-category}
Les arguments qui démontrent les lemmes
\ref{lemma-site-abelian-category} et \ref{lemma-embedding} valent aussi pour
les {\it catégories exactes} ; voir \cite[Appendix A]{Buhler} et
\cite[1.1.4]{BBD}. Nous les rappelons brièvement ici et ajouterons des détails
si le Champs Project en a besoin.

\medskip\noindent
Soit $\mathcal{A}$ une catégorie additive. Une paire {\it noyau-conoyau} est
une paire $(i, p)$ de morphismes de $\mathcal{A}$,
$i : A \to B$, $p : B \to C$, telle que $i$ soit le noyau de $p$ et $p$ le
conoyau de $i$. Étant donné un ensemble $\mathcal{E}$ de paires
noyau-conoyau, on dit que $i : A \to B$ est un
{\it monomorphisme admissible} si $(i, p) \in \mathcal{E}$ pour un certain
morphisme $p$. De même, on dit qu'un morphisme $p : B \to C$ est un
{\it épimorphisme admissible} si $(i, p) \in \mathcal{E}$ pour un certain
morphisme $i$. La paire $(\mathcal{A}, \mathcal{E})$ est appelée
{\it catégorie exacte} si les axiomes suivants sont satisfaits
\begin{enumerate}
\item $\mathcal{E}$ est stable par isomorphismes de paires noyau-conoyau ;
\item pour tout objet $A$, le morphisme $1_A$ est à la fois un épimorphisme
admissible et un monomorphisme admissible ;
\item les monomorphismes admissibles sont stables par composition ;
\item les épimorphismes admissibles sont stables par composition ;
\item la somme amalgamée d'un monomorphisme admissible $i : A \to B$ par tout
morphisme $A \to A'$ existe et le morphisme induit $i' : A' \to B'$ est un
monomorphisme admissible ;
\item le changement de base d'un épimorphisme admissible $p : B \to C$ par
tout morphisme $C' \to C$ existe et le morphisme induit $p' : B' \to C'$ est
un épimorphisme admissible.
\end{enumerate}
Étant donnée une telle structure, posons
$\mathcal{C} = (\mathcal{A}, \text{Cov})$, où les recouvrements, c'est-à-dire
les éléments de $\text{Cov}$, sont les épimorphismes admissibles. Les axiomes
ci-dessus entraînent immédiatement qu'il s'agit d'un site. Considérons le
foncteur
$$
F : \mathcal{A} \longrightarrow \textit{Ab}(\mathcal{C}), \quad
X \longmapsto h_X
$$
exactement comme au lemme \ref{lemma-embedding}. Ce foncteur est pleinement
fidèle, exact, et reflète l'exactitude. En outre, toute extension d'objets de
l'image essentielle de $F$ appartient à l'image essentielle de $F$.
\end{remark}






\section{Conditions AB de Grothendieck}
\label{section-grothendieck-conditions}

\noindent
Cette section et les quelques suivantes concernent surtout les « grandes »
catégories abéliennes, c'est-à-dire les catégories énumérées dans
Catégories, remarque \ref{categories-remark-big-categories}. L'exemple à
garder à l'esprit est la catégorie des faisceaux de modules sur un site annelé.

\medskip\noindent
Grothendieck a démontré l'existence d'injectifs dans une grande généralité dans
l'article \cite{Tohoku}. Il a utilisé les conditions suivantes pour distinguer
les catégories abéliennes dotées de propriétés particulières.

\begin{definition}
\label{definition-grothendieck-conditions}
Soit $\mathcal{A}$ une catégorie abélienne. Nous nommons quelques conditions :
\begin{enumerate}
\item[AB3] $\mathcal{A}$ possède des sommes directes ;
\item[AB4] $\mathcal{A}$ satisfait AB3 et les sommes directes sont exactes ;
\item[AB5] $\mathcal{A}$ satisfait AB3 et les colimites filtrantes sont
exactes.
\end{enumerate}
Voici les notions duales :
\begin{enumerate}
\item[AB3*] $\mathcal{A}$ possède des produits ;
\item[AB4*] $\mathcal{A}$ satisfait AB3* et les produits sont exacts ;
\item[AB5*] $\mathcal{A}$ satisfait AB3* et les limites cofiltrantes sont
exactes.
\end{enumerate}
On dit qu'un objet $U$ de $\mathcal{A}$ est un {\it générateur} si, pour tous
$N \subset M$, $N \not = M$, dans $\mathcal{A}$, il existe un morphisme
$U \to M$ qui ne se factorise pas par $N$. On dit que $\mathcal{A}$ est une
{\it catégorie abélienne de Grothendieck} si elle satisfait AB5 et possède un
générateur.
\end{definition}

\noindent
Discussion : une somme directe dans une catégorie abélienne est un coproduit.
Si une catégorie abélienne possède des sommes directes, c'est-à-dire satisfait
AB3, alors elle possède des colimites ; voir
Catégories, lemme \ref{categories-lemma-colimits-coproducts-coequalizers}. 
De même, si $\mathcal{A}$ satisfait AB3*, alors elle possède des limites ; voir
Catégories, lemme \ref{categories-lemma-limits-products-equalizers}. 
L'exactitude des sommes directes signifie ceci : étant donnés un ensemble
d'indices $I$ et des suites exactes courtes
$$
0 \to A_i \to B_i \to C_i \to 0,\quad i \in I
$$
dans $\mathcal{A}$, la suite
$$
0 \to
\bigoplus\nolimits_{i \in I} A_i \to
\bigoplus\nolimits_{i \in I} B_i \to
\bigoplus\nolimits_{i \in I} C_i \to 0
$$
est elle aussi exacte. Sans supposer AB4, on peut seulement affirmer en général
que cette suite est exacte à droite, c'est-à-dire que la prise des sommes
directes est un foncteur exact à droite lorsqu'elles existent. De même,
l'exactitude des colimites filtrantes signifie ceci : étant donnés un ensemble
filtrant $I$ et un système de suites exactes courtes
$$
0 \to A_i \to B_i \to C_i \to 0
$$
indexé par $I$ dans $\mathcal{A}$, la suite
$$
0 \to
\colim_{i \in I} A_i \to
\colim_{i \in I} B_i \to
\colim_{i \in I} C_i \to 0
$$
est elle aussi exacte. Sans supposer AB5, on peut seulement affirmer en général
que cette suite est exacte à droite, c'est-à-dire que la prise des colimites
est un foncteur exact à droite lorsqu'elles existent. Une explication analogue
vaut pour AB4* et AB5*.
\section{Objets injectifs dans les catégories de Grothendieck}
\label{section-grothendieck-categories}

\noindent
L'existence d'un générateur implique que, pour tout objet $M$ d'une
catégorie abélienne de Grothendieck $\mathcal{A}$, ses sous-objets
forment un ensemble. (Cela peut être faux pour une catégorie abélienne
``grande'' quelconque.)

\begin{lemma}
\label{lemma-set-of-subobjects}
Soient $\mathcal{A}$ une catégorie abélienne munie d'un générateur $U$ et
$X$ un objet de $\mathcal{A}$. Si $\kappa$ est le cardinal de
$\Mor(U, X)$, alors
\begin{enumerate}
\item Il n'existe aucune chaîne strictement croissante
(ni strictement décroissante) de sous-objets de $X$ indexée par
un cardinal strictement supérieur à $\kappa$.
\item Si $\alpha$ est un ordinal de cofinalité $> \kappa$,
toute suite croissante (ou décroissante) de sous-objets de $X$
indexée par $\alpha$ est stationnaire à partir d'un certain rang.
\item Le cardinal de l'ensemble des sous-objets de $X$
est $\leq 2^\kappa$.
\end{enumerate}
\end{lemma}

\begin{proof}
Pour (1), soit $\kappa' > \kappa$ un cardinal et supposons que
la famille $X_i$, $i \in \kappa'$, soit strictement croissante. Pour tout
$i$, choisissons alors $\phi_i \in \Mor(U, X)$ tel que $\phi_i$ se factorise
par $X_{i + 1}$, mais non par $X_i$. Les morphismes $\phi_i$ sont donc
deux à deux distincts, ce qui contredit la définition de $\kappa$.

\medskip\noindent
L'assertion (2) résulte de la définition de la cofinalité et de (1).

\medskip\noindent
Démonstration de (3). Pour tout sous-objet $Y \subset X$, définissons
$S_Y \in \mathcal{P}(\Mor(U, X))$ (l'ensemble des parties) par
$S_Y = \{\phi \in \Mor(U,X) : \phi\text{ se factorise par }Y\}$.
Alors $Y = Y'$ si et seulement si $S_Y = S_{Y'}$. Le cardinal de
l'ensemble des sous-objets est donc au plus celui de cet ensemble des parties.
\end{proof}

\noindent
D'après le lemme \ref{lemma-set-of-subobjects}, la définition suivante
a un sens.

\begin{definition}
\label{definition-size}
Soit $\mathcal{A}$ une catégorie abélienne de Grothendieck.
Soit $M$ un objet de $\mathcal{A}$.
La {\it taille} $|M|$ de $M$ est le cardinal de l'ensemble des sous-objets
de $M$.
\end{definition}

\begin{lemma}
\label{lemma-size-goes-down}
Soit $\mathcal{A}$ une catégorie abélienne de Grothendieck.
Si $0 \to M' \to M \to M'' \to 0$ est une suite exacte courte dans
$\mathcal{A}$, alors $|M'|, |M''| \leq |M|$.
\end{lemma}

\begin{proof}
Cela résulte immédiatement des définitions.
\end{proof}

\begin{lemma}
\label{lemma-set-iso-classes-bounded-size}
Soit $\mathcal{A}$ une catégorie abélienne de Grothendieck munie d'un
générateur $U$.
\begin{enumerate}
\item Si $|M| \leq \kappa$, alors $M$ est un quotient d'une somme directe
d'au plus $\kappa$ copies de $U$.
\item Pour tout cardinal $\kappa$, les classes d'isomorphisme des objets
$M$ tels que $|M| \leq \kappa$ forment un ensemble.
\end{enumerate}
\end{lemma}

\begin{proof}
Pour (1), choisissons, pour tout sous-objet propre $M' \subset M$, un morphisme
$\varphi_{M'} : U \to M$ dont l'image n'est pas contenue dans $M'$. Alors
$\bigoplus_{M' \subset M} \varphi_{M'} : \bigoplus_{M' \subset M} U \to M$
est surjectif. Il est clair que (1) implique (2).
\end{proof}

\begin{proposition}
\label{proposition-objects-are-small}
Soit $\mathcal{A}$ une catégorie abélienne de Grothendieck. Soit $M$ un
objet de $\mathcal{A}$. Posons $\kappa = |M|$.
Si $\alpha$ est un ordinal dont la cofinalité est strictement supérieure
à $\kappa$, alors $M$ est $\alpha$-petit relativement aux monomorphismes.
\end{proposition}

\begin{proof}
Comparons avec la proposition \ref{proposition-modules-are-small}.
Il suffit de montrer que l'application (\ref{equation-compare}) est surjective.
Soit $f : M \to \colim B_\beta$ une application.
Considérons les sous-objets $\{f^{-1}(B_\beta)\}$ de $M$, où $B_\beta$
est vu comme un sous-objet de la colimite $B = \bigcup_\beta B_\beta$.
Si l'un d'eux, disons $f^{-1}(B_\beta)$, est égal à $M$,
alors l'application se factorise par $B_\beta$.

\medskip\noindent
Puisque $\mathcal{A}$ vérifie AB5, on a
$$
\colim f^{-1}(B_\beta) = f^{-1}\left(\colim B_\beta\right) = M.
$$
Par hypothèse, il y a au plus $\kappa$ sous-objets distincts de $M$ parmi
les $f^{-1}(B_\alpha)$.
On peut donc trouver une partie $S \subset \alpha$ de cardinal au plus
$\kappa$ telle que, lorsque $\beta'$ parcourt $S$, les
$f^{-1}(B_{\beta'})$ parcourent \emph{tous} les $f^{-1}(B_\alpha)$.

\medskip\noindent
Toutefois, $S$ admet un majorant $\widetilde{\alpha} < \alpha$, puisque
$\alpha$ est de cofinalité strictement supérieure à $\kappa$. En particulier,
tous les $f^{-1}(B_{\beta'})$, $\beta' \in S$, sont contenus dans
$f^{-1}(B_{\widetilde{\alpha}})$.
Il s'ensuit que $f^{-1}(B_{\widetilde{\alpha}}) = M$.
L'application $f$ se factorise donc par $B_{\widetilde{\alpha}}$.
\end{proof}

\begin{lemma}
\label{lemma-characterize-injective}
\begin{slogan}
Pour vérifier qu'un objet est injectif, il suffit de vérifier la propriété
de relèvement pour les sous-objets d'un générateur.
\end{slogan}
Soit $\mathcal{A}$ une catégorie abélienne de Grothendieck munie d'un
générateur $U$.
Un objet $I$ de $\mathcal{A}$ est injectif si et seulement si, dans tout
diagramme commutatif
$$
\xymatrix{
M \ar[d] \ar[r] &  I \\
U \ar@{-->}[ru]
}
$$
où $M \subset U$ est un sous-objet, la flèche pointillée existe.
\end{lemma}

\begin{proof}
Voir le lemme \ref{lemma-criterion-baer} pour le cas des modules.
Soient un monomorphisme $A \subset B$ et un morphisme $\varphi : A \to I$.
Considérons l'ensemble $S$ des couples $(A', \varphi')$ formés de
sous-objets $A \subset A' \subset B$ et d'un morphisme $\varphi' : A' \to I$
prolongeant $\varphi$. Munissons cet ensemble de l'ordre partiel évident.
Soit $T \subset S$ une partie totalement ordonnée. Alors
$$
A' = \colim_{t \in T} A_t \xrightarrow{\colim_{t \in T} \varphi_t} I
$$
est un majorant. Le lemme de Zorn assure donc que $S$ possède un élément
maximal $(A', \varphi')$. Nous affirmons que $A' = B$. Sinon, choisissons
un morphisme $\psi : U \to B$ qui ne se factorise pas par $A'$. Posons
$N = A' \cap \psi(U)$ et $M = \psi^{-1}(N)$. L'application
$$
M \to N \to A' \xrightarrow{\varphi'} I
$$
se prolonge en un morphisme $\chi : U \to I$. Comme
$\chi|_{\Ker(\psi)} = 0$, le morphisme $\chi$ se factorise sous la forme
$$
U \to \Im(\psi) \xrightarrow{\varphi''} I
$$
Puisque $\varphi'$ et $\varphi''$ coïncident sur
$N = A' \cap \Im(\psi)$, ils définissent ensemble un morphisme
$A' + \Im(\psi) \to I$, ce qui contredit la maximalité supposée de $A'$.
\end{proof}

\begin{theorem}
\label{theorem-injective-embedding-grothendieck}
Soit $\mathcal{A}$ une catégorie abélienne de Grothendieck.
Alors $\mathcal{A}$ admet des plongements injectifs fonctoriels.
\end{theorem}

\begin{proof}
Comparons avec la démonstration du
théorème \ref{theorem-baer-grothendieck}.
Choisissons un générateur $U$ de $\mathcal{A}$. Pour un objet $M$, on définit
$\mathbf{M}(M)$ par le diagramme cocartésien suivant :
$$
\xymatrix{
\bigoplus_{N \subset U}
\bigoplus_{\varphi \in \Hom(N, M)}
N \ar[r] \ar[d] & M \ar[d] \\
\bigoplus_{N \subset U}
\bigoplus_{\varphi \in \Hom(N, M)}
U \ar[r] &  \mathbf{M}(M).
}
$$
Notons que $M \mapsto \mathbf{M}(M)$ est un foncteur et qu'il existe des
applications injectives fonctorielles $M \to \mathbf{M}(M)$. On définit par
récurrence transfinie des foncteurs $\mathbf{M}_\alpha(M)$ pour tout ordinal
$\alpha$. On pose $\mathbf{M}_0(M) = M$. Une fois $\mathbf{M}_\alpha(M)$
défini, on pose
$\mathbf{M}_{\alpha + 1}(M) = \mathbf{M}(\mathbf{M}_\alpha(M))$.
Pour un ordinal limite $\beta$, on pose
$$
\mathbf{M}_\beta(M) = \colim_{\alpha < \beta} \mathbf{M}_\alpha(M).
$$
Enfin, choisissons un ordinal $\alpha$ dont la cofinalité est strictement
supérieure à $|U|$. Un tel ordinal existe d'après
Ensembles, proposition \ref{sets-proposition-exist-ordinals-large-cofinality}.
Nous affirmons que $M \to \mathbf{M}_\alpha(M)$ est le plongement injectif
fonctoriel recherché. En effet, si $N \subset U$ est un sous-objet et si
$\varphi : N \to \mathbf{M}_\alpha(M)$ est un morphisme, alors
$\varphi$ se factorise par $\mathbf{M}_{\alpha'}(M)$ pour un certain
$\alpha' < \alpha$, d'après la
proposition \ref{proposition-objects-are-small}.
Par construction de $\mathbf{M}(-)$, le morphisme $\varphi$ se prolonge en
un morphisme de $U$ dans $\mathbf{M}_{\alpha' + 1}(M)$, donc dans
$\mathbf{M}_\alpha(M)$. Le
lemme \ref{lemma-characterize-injective}
montre alors que $\mathbf{M}_\alpha(M)$ est injectif.
\end{proof}













\section{Complexes K-injectifs dans les catégories de Grothendieck}
\label{section-K-injective}

\noindent
Le contenu de cette section provient de l'article \cite{serpe} de Serp\'e.
Cet article généralise certains résultats de Spaltenstein
\cite{Spaltenstein} à des catégories abéliennes de Grothendieck quelconques.
Notre lemme \ref{lemma-characterize-K-injective} n'apparaît
qu'implicitement dans l'article de Serp\'e. Notre démarche consiste à imiter
la démonstration de Grothendieck du
théorème \ref{theorem-injective-embedding-grothendieck}.

\begin{lemma}
\label{lemma-surjection-bounded-size}
Soit $\mathcal{A}$ une catégorie abélienne de Grothendieck munie d'un
générateur $U$. Soit $c$ la fonction sur les cardinaux définie par
$c(\kappa) = |\bigoplus_{\alpha \in \kappa} U|$. Si $\pi : M \to N$ est
surjectif, il existe un sous-objet $M' \subset M$ qui se surjecte sur $N$
et tel que $|M'| \leq c(|N|)$.
\end{lemma}

\begin{proof}
Pour tout sous-objet propre $N' \subset N$, choisissons un morphisme
$\varphi_{N'} : U \to M$ tel que $U \to M \to N$ ne se factorise pas
par $N'$. Posons
$$
M' = \Im\left(
\bigoplus\nolimits_{N' \subset N} \varphi_{N'} :
\bigoplus\nolimits_{N' \subset N} U \longrightarrow M\right)
$$
Alors $M'$ convient.
\end{proof}

\begin{lemma}
\label{lemma-acyclic-quotient-complexes-bounded-size}
Soit $\mathcal{A}$ une catégorie abélienne de Grothendieck. Il existe un
cardinal $\kappa$ tel que, pour tout complexe acyclique $M^\bullet$, on ait :
\begin{enumerate}
\item si $M^\bullet$ est non nul, il existe un sous-complexe non nul
$N^\bullet$, borné supérieurement et acyclique, tel que
$|N^n| \leq \kappa$ ;
\item il existe une surjection de complexes
$$
\bigoplus\nolimits_{i \in I} M_i^\bullet \longrightarrow M^\bullet
$$
où $M_i^\bullet$ est borné supérieurement, acyclique, et
$|M_i^n| \leq \kappa$.
\end{enumerate}
\end{lemma}

\begin{proof}
Choisissons un générateur $U$ de $\mathcal{A}$. Notons $c$ la fonction du
lemme \ref{lemma-surjection-bounded-size}.
Posons $\kappa = \sup \{c^n(|U|), n = 1, 2, 3, \ldots\}$.
Soient $n \in \mathbf{Z}$ et un morphisme $\psi : U \to M^n$.
Pour démontrer (1) et (2), il suffit d'établir qu'il existe un sous-complexe
$N^\bullet \subset M^\bullet$, borné supérieurement et acyclique, tel que
$|N^m| \leq \kappa$ et que $\psi$ se factorise par $N^n$.
À cette fin, posons $N^n = \Im(\psi)$,
$N^{n + 1} = \Im(U \to M^n \to M^{n + 1})$, et $N^m = 0$ pour
$m \geq n + 2$.
Supposons que, pour un certain $k \leq n$, on ait construit
$N^m \subset M^m$ pour tout $n > m \geq k$, de telle sorte que
\begin{enumerate}
\item $\text{d}(N^m) \subset N^{m + 1}$ pour tout $m \geq k$ ;
\item $\Im(N^{m - 1} \to N^m) = \Ker(N^m \to N^{m + 1})$ pour
tout $m \geq k + 1$ ;
\item $|N^m| \leq c^{\max\{n - m, 0\}}(|U|)$ pour tout $m \geq k$.
\end{enumerate}
Comme $M^\bullet$ est acyclique, le sous-objet
$\text{d}^{-1}(\Ker(N^k \to N^{k + 1})) \subset M^{k - 1}$ se surjecte sur
$\Ker(N^k \to N^{k + 1})$. Le
lemme \ref{lemma-surjection-bounded-size} permet donc de choisir
$N^{k - 1} \subset M^{k - 1}$ se surjectant sur
$\Ker(N^k \to N^{k + 1})$ et vérifiant
$|N^{k - 1}| \leq c^{n - k + 1}(|U|)$.
On conclut par récurrence sur $k$.
\end{proof}

\begin{lemma}
\label{lemma-characterize-K-injective}
Soit $\mathcal{A}$ une catégorie abélienne de Grothendieck.
Soit $\kappa$ un cardinal comme dans le
lemme \ref{lemma-acyclic-quotient-complexes-bounded-size}.
Supposons que $I^\bullet$ soit un complexe tel que
\begin{enumerate}
\item chaque $I^j$ soit injectif ;
\item pour tout complexe acyclique $M^\bullet$, borné supérieurement,
tel que $|M^n| \leq \kappa$, on ait
$\Hom_{K(\mathcal{A})}(M^\bullet, I^\bullet) = 0$.
\end{enumerate}
Alors $I^\bullet$ est un complexe $K$-injectif.
\end{lemma}

\begin{proof}
Soit $M^\bullet$ un complexe acyclique. Nous allons construire, par
récurrence sur l'ordinal $\alpha$, un sous-complexe acyclique
$K_\alpha^\bullet \subset M^\bullet$ de la manière suivante.
Pour $\alpha = 0$, posons $K_0^\bullet = 0$. Pour $\alpha > 0$,
procédons comme suit :
\begin{enumerate}
\item Si $\alpha = \beta + 1$ et $K_\beta^\bullet = M^\bullet$,
on choisit $K_\alpha^\bullet = K_\beta^\bullet$.
\item Si $\alpha = \beta + 1$ et $K_\beta^\bullet \not = M^\bullet$,
alors $M^\bullet/K_\beta^\bullet$ est un complexe acyclique non nul.
On choisit un sous-complexe
$N_\alpha^\bullet \subset M^\bullet/K_\beta^\bullet$ comme dans le
lemme \ref{lemma-acyclic-quotient-complexes-bounded-size}.
Enfin, on prend pour $K_\alpha^\bullet \subset M^\bullet$
l'image inverse de $N_\alpha^\bullet$.
\item Si $\alpha$ est un ordinal limite, on pose
$K_\alpha^\bullet = \colim_{\beta < \alpha} K_\beta^\bullet$.
\end{enumerate}
Il est clair que $M^\bullet = K_\alpha^\bullet$ pour un ordinal $\alpha$
assez grand. Nous allons démontrer que
$$
\Hom_{K(\mathcal{A})}(K_\alpha^\bullet, I^\bullet)
$$
est nul par récurrence transfinie sur $\alpha$. C'est vrai pour $\alpha = 0$,
car $K_0^\bullet$ est nul. Supposons le résultat vrai pour $\beta$ et
$\alpha = \beta + 1$. Dans le cas (1) de la liste ci-dessus, le résultat est
clair. Dans le cas (2), on dispose d'une suite exacte courte de complexes
$$
0 \to K_\beta^\bullet \to K_\alpha^\bullet \to N_\alpha^\bullet \to 0
$$
Puisque chaque composante de $I^\bullet$ est injective, on obtient une suite
exacte
$$
\Hom_{K(\mathcal{A})}(K_\beta^\bullet, I^\bullet) \leftarrow
\Hom_{K(\mathcal{A})}(K_\alpha^\bullet, I^\bullet) \leftarrow
\Hom_{K(\mathcal{A})}(N_\alpha^\bullet, I^\bullet)
$$
Par récurrence, le terme de gauche est nul et, par hypothèse sur
$I^\bullet$, celui de droite est nul. Le groupe du milieu est donc lui aussi
nul. Supposons enfin que $\alpha$ soit un ordinal limite. Alors
$$
\Hom^\bullet(K_\alpha^\bullet, I^\bullet) =
\lim_{\beta < \alpha} \Hom^\bullet(K_\beta^\bullet, I^\bullet)
$$
avec les notations de
Compléments d'algèbre, section \ref{more-algebra-section-hom-complexes}.
Ces complexes calculent les morphismes dans $K(\mathcal{A})$ d'après
Compléments d'algèbre, équation
(\ref{more-algebra-equation-cohomology-hom-complex}).
Les applications de transition du système sont surjectives, puisque $I^j$
est injectif pour tout $j$. En outre, pour un ordinal limite $\alpha$, la
limite est égale à la valeur (voir la formule affichée ci-dessus). On peut donc
conclure à l'aide de
Homologie, lemme \ref{homology-lemma-ML-over-ordinals}.
\end{proof}

\begin{lemma}
\label{lemma-functorial-homotopies}
Soit $\mathcal{A}$ une catégorie abélienne de Grothendieck.
Soit $(K_i^\bullet)_{i \in I}$ un ensemble de complexes acycliques.
Il existe un foncteur $M^\bullet \mapsto \mathbf{M}^\bullet(M^\bullet)$
et une transformation naturelle
$j_{M^\bullet} : M^\bullet \to \mathbf{M}^\bullet(M^\bullet)$
tels que
\begin{enumerate}
\item $j_{M^\bullet}$ soit un quasi-isomorphisme injectif (terme à terme) ;
\item pour tout $i \in I$ et tout $w : K_i^\bullet \to M^\bullet$,
le morphisme $j_{M^\bullet} \circ w$ soit homotope à zéro.
\end{enumerate}
\end{lemma}

\begin{proof}
Pour tout $i \in I$, choisissons une application de complexes injective
terme à terme $K_i^\bullet \to L_i^\bullet$, homotope à zéro, où
$L_i^\bullet$ est quasi-isomorphe à zéro. On peut par exemple prendre pour
$L_i^\bullet$ le cône de l'identité de $K_i^\bullet$.
On définit $\mathbf{M}^\bullet(M^\bullet)$ par le diagramme cocartésien suivant :
$$
\xymatrix{
\bigoplus_{i \in I}
\bigoplus_{w : K_i^\bullet \to M^\bullet}
K_i^\bullet \ar[r] \ar[d] & M^\bullet \ar[d] \\
\bigoplus_{i \in I}
\bigoplus_{w : K_i^\bullet \to M^\bullet}
L_i^\bullet \ar[r] &  \mathbf{M}^\bullet(M^\bullet).
}
$$
Alors $M^\bullet \mapsto \mathbf{M}^\bullet(M^\bullet)$ est un foncteur. La flèche
verticale de droite définit l'application injective fonctorielle
$j_{M^\bullet}$. Le conoyau de $j_{M^\bullet}$ est isomorphe à la somme directe
des conoyaux des
applications $K_i^\bullet \to L_i^\bullet$, et est donc acyclique. Ainsi,
$j_{M^\bullet}$ est un quasi-isomorphisme. L'assertion (2) est vraie par
construction.
\end{proof}

\begin{lemma}
\label{lemma-functorial-injective}
Soit $\mathcal{A}$ une catégorie abélienne de Grothendieck.
Il existe un foncteur $M^\bullet \mapsto \mathbf{N}^\bullet(M^\bullet)$
et une transformation naturelle
$j_{M^\bullet} : M^\bullet \to \mathbf{N}^\bullet(M^\bullet)$
tels que
\begin{enumerate}
\item $j_{M^\bullet}$ soit un quasi-isomorphisme injectif (terme à terme) ;
\item pour tout $n \in \mathbf{Z}$, l'application
$M^n \to \mathbf{N}^n(M^\bullet)$ se factorise par un sous-objet
$I^n \subset \mathbf{N}^n(M^\bullet)$, où $I^n$ est un objet injectif de
$\mathcal{A}$.
\end{enumerate}
\end{lemma}

\begin{proof}
Choisissons des plongements injectifs fonctoriels $i_M : M \to I(M)$, voir le
théorème \ref{theorem-injective-embedding-grothendieck}.
Pour tout complexe $M^\bullet$, notons $J^\bullet(M^\bullet)$ le complexe
dont les termes sont $J^n(M^\bullet) = I(M^n) \oplus I(M^{n + 1})$ et dont
la différentielle est
$$
d_{J^\bullet(M^\bullet)} =
\left(
\begin{matrix}
0 & 1 \\
0 & 0
\end{matrix}
\right)
$$
Il existe une application injective canonique de complexes
$u_{M^\bullet} : M^\bullet \to J^\bullet(M^\bullet)$, obtenue en envoyant
$M^n$ dans $I(M^n) \oplus I(M^{n + 1})$ au moyen des applications
$i_{M^n} : M^n \to I(M^n)$ et
$i_{M^{n + 1}} \circ d : M^n \to M^{n + 1} \to I(M^{n + 1})$. On dispose donc
d'une suite exacte courte de complexes
$$
0 \to M^\bullet \xrightarrow{u_{M^\bullet}}
J^\bullet(M^\bullet) \xrightarrow{v_{M^\bullet}}
Q^\bullet(M^\bullet) \to 0
$$
fonctorielle en $M^\bullet$. Posons
$$
\mathbf{N}^\bullet(M^\bullet) = C(v_{M^\bullet})^\bullet[-1].
$$
Remarquons que
$$
\mathbf{N}^n(M^\bullet) = Q^{n - 1}(M^\bullet) \oplus J^n(M^\bullet)
$$
dont la différentielle est
$$
\left(
\begin{matrix}
- d^{n - 1}_{Q^\bullet(M^\bullet)} & - v^n_{M^\bullet} \\
0 & d^n_{J^\bullet(M^\bullet)}
\end{matrix}
\right)
$$
Il existe donc une application de complexes
$j_{M^\bullet} : M^\bullet \to \mathbf{N}^\bullet(M^\bullet)$ induite par $u$.
Elle est injective et se factorise par un sous-objet injectif par construction.
L'application $j_{M^\bullet}$ est un quasi-isomorphisme, comme le montre la
suite exacte longue de cohomologie associée aux suites exactes courtes de
complexes ci-dessus.
\end{proof}

\begin{theorem}
\label{theorem-K-injective-embedding-grothendieck}
\begin{slogan}
Existence de complexes K-injectifs dans les catégories abéliennes de
Grothendieck.
\end{slogan}
Soit $\mathcal{A}$ une catégorie abélienne de Grothendieck.
Pour tout complexe $M^\bullet$, il existe un quasi-isomorphisme
$M^\bullet \to I^\bullet$ tel que $M^n \to I^n$ soit injectif, que $I^n$
soit un objet injectif de $\mathcal{A}$ pour tout $n$, et que $I^\bullet$
soit un complexe K-injectif. En outre, cette construction est fonctorielle
en $M^\bullet$.
\end{theorem}

\begin{proof}
Comparons avec les démonstrations des
théorèmes \ref{theorem-baer-grothendieck} et
\ref{theorem-injective-embedding-grothendieck}.
Choisissons un cardinal $\kappa$ comme dans les
lemmes \ref{lemma-acyclic-quotient-complexes-bounded-size} et
\ref{lemma-characterize-K-injective}.
Choisissons un ensemble $(K_i^\bullet)_{i \in I}$ de complexes acycliques
bornés supérieurement tel que tout complexe acyclique $K^\bullet$, borné
supérieurement et vérifiant $|K^n| \leq \kappa$, soit isomorphe à
$K_i^\bullet$ pour un certain $i \in I$. Cela est possible d'après le
lemme \ref{lemma-set-iso-classes-bounded-size}.
Notons $\mathbf{M}^\bullet(-)$ le foncteur construit dans le
lemme \ref{lemma-functorial-homotopies}, et
$\mathbf{N}^\bullet(-)$ celui construit dans le
lemme \ref{lemma-functorial-injective}.
Ces deux foncteurs sont munis de transformations injectives
$\text{id} \to \mathbf{M}$ et $\text{id} \to \mathbf{N}$.

\medskip\noindent
Par récurrence transfinie, on définit une suite de foncteurs
$\mathbf{T}_\alpha(-)$ et les transformations correspondantes
$\text{id} \to \mathbf{T}_\alpha$. On pose
$\mathbf{T}_0(M^\bullet) = M^\bullet$. Si $\mathbf{T}_\alpha$ est défini,
on pose
$$
\mathbf{T}_{\alpha + 1}(M^\bullet) =
\mathbf{N}^\bullet(\mathbf{M}^\bullet(\mathbf{T}_\alpha(M^\bullet)))
$$
Si $\beta$ est un ordinal limite, on pose
$$
\mathbf{T}_\beta(M^\bullet) =
\colim_{\alpha < \beta} \mathbf{T}_\alpha(M^\bullet)
$$
Les applications de transition du système sont des quasi-isomorphismes
injectifs. D'après AB5, la colimite est encore quasi-isomorphe à $M^\bullet$.
Nous affirmons que $M^\bullet \to \mathbf{T}_\alpha(M^\bullet)$ convient si
la cofinalité de $\alpha$ est strictement supérieure à
$\max(\kappa, |U|)$, où $U$ est un générateur de $\mathcal{A}$.
Il suffit en effet de vérifier les conditions (1) et (2) du
lemme \ref{lemma-characterize-K-injective}.

\medskip\noindent
Pour (1), utilisons le critère du
lemme \ref{lemma-characterize-injective}.
Supposons que $M \subset U$ et que
$\varphi : M \to \mathbf{T}^n_\alpha(M^\bullet)$ soit un morphisme pour un
certain $n \in \mathbf{Z}$. La
proposition \ref{proposition-objects-are-small} montre que $\varphi$ se
factorise par $\mathbf{T}^n_{\alpha'}(M^\bullet)$ pour un certain
$\alpha' < \alpha$. En particulier, par construction du foncteur
$\mathbf{N}^\bullet(-)$, le morphisme $\varphi$ se factorise par un objet
injectif de $\mathcal{A}$ ; cela montre que $\varphi$ se relève en un morphisme
défini sur $U$.

\medskip\noindent
Pour (2), soit $w : K^\bullet \to \mathbf{T}_\alpha(M^\bullet)$ un morphisme
de complexes, où $K^\bullet$ est un complexe acyclique borné supérieurement
tel que $|K^n| \leq \kappa$. Alors $K^\bullet \cong K_i^\bullet$ pour un
certain $i \in I$. De nouveau, la
proposition \ref{proposition-objects-are-small} montre que $w$ se factorise
par $\mathbf{T}_{\alpha'}(M^\bullet)$ pour un certain $\alpha' < \alpha$.
Par construction du foncteur $\mathbf{M}^\bullet(-)$, il s'ensuit que $w$
est homotope à zéro. Cela achève la démonstration.
\end{proof}







\section{Remarques complémentaires sur les catégories abéliennes de
Grothendieck}
\label{section-additional-Grothendieck}

\noindent
Nous rassemblons dans cette section quelques résultats classiques sur les
catégories abéliennes de Grothendieck.

\begin{lemma}
\label{lemma-grothendieck-brown}
Soit $\mathcal{A}$ une catégorie abélienne de Grothendieck.
Soit $F : \mathcal{A}^{opp} \to \textit{Ensembles}$ un foncteur.
Alors $F$ est représentable si et seulement si $F$ transforme les colimites
en limites, c'est-à-dire si
$$
F(\colim_i N_i) = \lim F(N_i)
$$
pour tout diagramme $\mathcal{I} \to \mathcal{A}$, $i \in \mathcal{I}$.
\end{lemma}

\begin{proof}
Si $F$ est représentable, il transforme les colimites en limites par
définition des colimites.

\medskip\noindent
Supposons que $F$ transforme les colimites en limites. Alors
$F(M \oplus N) = F(M) \times F(N)$, ce qui permet de munir $F(M)$ d'une
structure de groupe. On obtient ainsi un foncteur
$F : \mathcal{A}^{opp} \to \textit{Ab}$ additif et exact à droite au sens
contravariant, c'est-à-dire qu'il transforme une suite exacte courte
$0 \to K \to L \to M \to 0$ en une suite exacte
$F(K) \leftarrow F(L) \leftarrow F(M) \leftarrow 0$ (comparer avec
Homologie, section \ref{homology-section-functors}).

\medskip\noindent
Soit $U$ un générateur de $\mathcal{A}$. Posons
$A = \bigoplus_{s \in F(U)} U$ et
$s_{univ} = (s)_{s \in F(U)} \in F(A) = \prod_{s \in F(U)} F(U)$.
Soit $A' \subset A$ le plus grand sous-objet tel que la restriction de
$s_{univ}$ à $A'$ soit nulle. Il existe parce que $\mathcal{A}$ est une
catégorie de Grothendieck et que $F$ transforme les colimites en limites.
Puisque $F$ transforme les colimites en limites, il existe aussi un unique
élément $\overline{s}_{univ} \in F(A/A')$ qui s'envoie sur $s_{univ}$ dans
$F(A)$.
Nous affirmons que $A/A'$ représente $F$ ; autrement dit, l'application de
Yoneda
$$
\overline{s}_{univ} : h_{A/A'} \longrightarrow F
$$
est un isomorphisme. Soient $M \in \Ob(\mathcal{A})$ et $s \in F(M)$.
Considérons la surjection
$$
c_M :
A_M = \bigoplus\nolimits_{\varphi \in \Hom_\mathcal{A}(U, M)} U
\longrightarrow
M.
$$
Elle donne $F(c_M)(s) = (s_\varphi) \in \prod_\varphi F(U)$.
Considérons l'application
$$
\psi :
A_M = \bigoplus\nolimits_{\varphi \in \Hom_\mathcal{A}(U, M)} U
\longrightarrow
\bigoplus\nolimits_{s \in F(U)} U = A
$$
qui envoie le facteur correspondant à $\varphi$ sur celui correspondant à
$s_\varphi$ par l'application identique de $U$. Par construction, $s_{univ}$
a pour image $(s_\varphi)_\varphi$. Autrement dit, le carré de droite du
diagramme
$$
\xymatrix{
A' \ar[r] &
A \ar@{..>}[r]_{s_{univ}} & F \\
K \ar[r] \ar[u]^{?} & A_M \ar[u]^\psi \ar[r] &
M \ar@{..>}[u]_s
}
$$
est commutatif. Posons $K = \Ker(A_M \to M)$. Comme la restriction de $s$
à $K$ est nulle, on a $\psi(K) \subset A'$ par définition de $A'$. Il existe
donc un morphisme induit $M \to A/A'$. Cette construction fournit un inverse
de l'application $h_{A/A'}(M) \to F(M)$ (nous omettons les détails).
\end{proof}

\begin{lemma}
\label{lemma-grothendieck-products}
Toute catégorie abélienne de Grothendieck vérifie Ab3*.
\end{lemma}

\begin{proof}
Soit $M_i$, $i \in I$, une famille d'objets de $\mathcal{A}$ indexée par un
ensemble $I$. Le foncteur $F = \prod_{i \in I} h_{M_i}$ transforme les
colimites en limites. Le
lemme \ref{lemma-grothendieck-brown} s'applique donc.
\end{proof}

\begin{remark}
\label{remark-existence-D}
Dans le chapitre consacré aux catégories dérivées, nous travaillons
systématiquement avec des catégories abéliennes ``petites'' (conformément à
la convention du Champs project). Pour une catégorie abélienne ``grande''
$\mathcal{A}$, l'existence de la catégorie dérivée $D(\mathcal{A})$ n'est pas
évidente, car il n'est pas clair que les morphismes de la catégorie dérivée
forment des ensembles. En général, c'est faux ; voir
Exemples, lemme \ref{examples-lemma-big-abelian-category}.
En revanche, si $\mathcal{A}$ est une catégorie abélienne de Grothendieck et si
$K^\bullet, L^\bullet$ sont dans $K(\mathcal{A})$, alors le
théorème \ref{theorem-K-injective-embedding-grothendieck} donne un
quasi-isomorphisme $L^\bullet \to I^\bullet$ vers un complexe K-injectif
$I^\bullet$, et Catégories dérivées,
lemme \ref{derived-lemma-K-injective}, montre que
$$
\Hom_{D(\mathcal{A})}(K^\bullet, L^\bullet) =
\Hom_{K(\mathcal{A})}(K^\bullet, I^\bullet)
$$
qui est un ensemble. La catégorie des modules sur un anneau et, plus
généralement, celle des faisceaux de modules sur un site annelé sont des
exemples de catégories abéliennes de Grothendieck.
\end{remark}

\begin{lemma}
\label{lemma-derived-products}
Soit $\mathcal{A}$ une catégorie abélienne de Grothendieck. Alors :
\begin{enumerate}
\item $D(\mathcal{A})$ admet des sommes directes et des produits ;
\item les sommes directes s'obtiennent en prenant terme à terme les sommes
directes de complexes quelconques ;
\item les produits s'obtiennent en prenant terme à terme les produits de
complexes K-injectifs.
\end{enumerate}
\end{lemma}

\begin{proof}
Soit $K^\bullet_i$, $i \in I$, une famille d'objets de $D(\mathcal{A})$
indexée par un ensemble $I$. Nous affirmons que la somme directe terme à
terme $\bigoplus_{i \in I} K^\bullet_i$ est une somme directe dans
$D(\mathcal{A})$. Soit en effet $I^\bullet$ un complexe K-injectif. On a
\begin{align*}
\Hom_{D(\mathcal{A})}(\bigoplus\nolimits_{i \in I} K^\bullet_i, I^\bullet)
& =
\Hom_{K(\mathcal{A})}(\bigoplus\nolimits_{i \in I} K^\bullet_i, I^\bullet) \\
& =
\prod\nolimits_{i \in I} \Hom_{K(\mathcal{A})}(K^\bullet_i, I^\bullet) \\
& =
\prod\nolimits_{i \in I} \Hom_{D(\mathcal{A})}(K^\bullet_i, I^\bullet)
\end{align*}
comme souhaité. Cela suffit, puisque tout complexe peut être représenté par
un complexe K-injectif d'après le
théorème \ref{theorem-K-injective-embedding-grothendieck}.
Pour construire le produit, choisissons une résolution K-injective
$K_i^\bullet \to I_i^\bullet$ pour chaque $i$. Nous affirmons alors que
$\prod_{i \in I} I_i^\bullet$ est un produit dans $D(\mathcal{A})$.
Cela résulte de Catégories dérivées,
lemme \ref{derived-lemma-product-K-injective}.
\end{proof}

\begin{remark}
\label{remark-direct-sum-product-derived}
Soit $R$ un anneau. Supposons que $M_n$, $n \in \mathbf{Z}$, soient des
$R$-modules. Notons $E_n = M_n[-n] \in D(R)$. Nous affirmons que
$E = \bigoplus M_n[-n]$ est {\it à la fois} la somme directe et le produit des
objets $E_n$ dans $D(R)$. Pour la somme directe, il suffit de consulter la
démonstration du lemme \ref{lemma-derived-products}.
Pour le produit direct, choisissons des résolutions injectives
$M_n \to I_n^\bullet$. La démonstration du
lemme \ref{lemma-derived-products} donne
$$
\prod E_n = \prod I_n^\bullet[-n]
$$
dans $D(R)$. Comme les produits sont exacts dans $\text{Mod}_R$,
$\prod I_n^\bullet[-n]$ est quasi-isomorphe à $E$. Le même raisonnement vaut
plus généralement dans $D(\mathcal{A})$, lorsque $\mathcal{A}$ est une
catégorie abélienne de Grothendieck vérifiant Ab4*.
\end{remark}

\begin{lemma}
\label{lemma-RF-commutes-with-Rlim}
Soit $F : \mathcal{A} \to \mathcal{B}$ un foncteur additif entre catégories
abéliennes. Supposons que
\begin{enumerate}
\item $\mathcal{A}$ soit une catégorie abélienne de Grothendieck ;
\item les produits dénombrables soient exacts dans $\mathcal{B}$ ;
\item $F$ commute aux produits dénombrables.
\end{enumerate}
Alors $RF : D(\mathcal{A}) \to D(\mathcal{B})$ commute aux limites dérivées.
\end{lemma}

\begin{proof}
Remarquons que $RF$ existe, car $\mathcal{A}$ possède suffisamment de
complexes K-injectifs (théorème
\ref{theorem-K-injective-embedding-grothendieck} et
Catégories dérivées, lemme \ref{derived-lemma-K-injective-defined}).
L'énoncé signifie que, si $K = R\lim K_n$, alors
$RF(K) = R\lim RF(K_n)$. Pour les notations, voir
Catégories dérivées, définition \ref{derived-definition-derived-limit}.
Comme $RF$ est un foncteur exact entre catégories triangulées, il suffit de
vérifier que $RF$ commute aux produits dénombrables d'objets de
$D(\mathcal{A})$. La démonstration du
lemme \ref{lemma-derived-products} montre que les produits dans
$D(\mathcal{A})$ se calculent en prenant les produits de complexes
K-injectifs, et qu'un produit de complexes K-injectifs est K-injectif.
En outre, Catégories dérivées,
lemme \ref{derived-lemma-products}, montre que les produits dans
$D(\mathcal{B})$ se calculent terme à terme.
Or $RF$ se calcule en appliquant $F$ à un représentant K-injectif, et nous
avons supposé que $F$ commute aux produits dénombrables ; le résultat en
découle.
\end{proof}
\noindent
Le lemme suivant constitue une sorte de généralisation de
l'existence des résolutions de Cartan--Eilenberg
(Catégories dérivées, section \ref{derived-section-cartan-eilenberg}).

\begin{lemma}
\label{lemma-K-injective-embedding-filtration}
Soit $\mathcal{A}$ une catégorie abélienne de Grothendieck.
Soit $K^\bullet$ un complexe filtré de $\mathcal{A}$, voir
Homologie, définition \ref{homology-definition-filtered-complex}.
Il existe alors un morphisme $j : K^\bullet \to J^\bullet$
de complexes filtrés de $\mathcal{A}$ tel que
\begin{enumerate}
\item $J^n$, $F^pJ^n$, $J^n/F^pJ^n$ et $F^pJ^n/F^{p'}J^n$ sont des
objets injectifs de $\mathcal{A}$,
\item $J^\bullet$, $F^pJ^\bullet$, $J^\bullet/F^pJ^\bullet$ et
$F^pJ^\bullet/F^{p'}J^\bullet$ sont des complexes K-injectifs,
\item $j$ induit les quasi-isomorphismes
$K^\bullet \to J^\bullet$,
$F^pK^\bullet \to F^pJ^\bullet$,
$K^\bullet/F^pK^\bullet \to J^\bullet/F^pJ^\bullet$, et
$F^pK^\bullet/F^{p'}K^\bullet \to F^pJ^\bullet/F^{p'}J^\bullet$.
\end{enumerate}
\end{lemma}

\begin{proof}
D'après le théorème \ref{theorem-K-injective-embedding-grothendieck},
on obtient des quasi-isomorphismes $i : K^\bullet \to I^\bullet$ et
$i^p : F^pK^\bullet \to I^{p, \bullet}$ ainsi que les diagrammes commutatifs
$$
\vcenter{
\xymatrix{
K^\bullet \ar[d]_i & F^pK^\bullet \ar[l] \ar[d]_{i^p} \\
I^\bullet & I^{p, \bullet} \ar[l]_{\alpha^p}
}
}
\quad\text{et}\quad
\vcenter{
\xymatrix{
F^{p'}K^\bullet \ar[d]_{i^{p'}} &
F^pK^\bullet \ar[l] \ar[d]_{i^p} \\
I^{p', \bullet} &
I^{p, \bullet} \ar[l]_{\alpha^{p p'}}
}
}
\quad\text{pour }p' \leq p
$$
de sorte que $\alpha^{p'} \circ \alpha^{p p'} = \alpha^p$
et $\alpha^{p'p''} \circ \alpha^{pp'} = \alpha^{pp''}$.
Le problème est que les applications $\alpha^p : I^{p, \bullet} \to I^\bullet$
ne sont pas nécessairement injectives. Pour chaque $p$, choisissons une
injection $t^p : I^{p, \bullet} \to J^{p, \bullet}$ dans un complexe
K-injectif acyclique $J^{p, \bullet}$ dont les termes sont des objets
injectifs de $\mathcal{A}$ (on l'envoie d'abord dans le cône de l'identité,
puis on applique le théorème).
Choisissons un morphisme de complexes $s^p : I^\bullet \to J^{p, \bullet}$
tel que le diagramme suivant soit commutatif
$$
\xymatrix{
K^\bullet \ar[d]_i & F^pK^\bullet \ar[l] \ar[d]_{i^p} \\
I^\bullet \ar[rd]_{s^p} & I^{p, \bullet} \ar[d]^{t^p} \\
& J^{p, \bullet}
}
$$
C'est possible : la composée $F^pK^\bullet \to J^{p, \bullet}$
est homotope à zéro, car $J^{p, \bullet}$ est acyclique et K-injectif
(Catégories dérivées, lemme \ref{derived-lemma-K-injective}).
Puisque les objets $J^{p, n - 1}$ sont injectifs et que
$F^pK^n \to K^n \to I^n$ sont des morphismes injectifs, on peut
prolonger les applications $F^pK^n \to J^{p, n - 1}$ qui définissent
l'homotopie en une application $h^n : I^n \to J^{p, n - 1}$.
On pose alors $s^p$ égal à $h \circ \text{d} + \text{d} \circ h$.
(Attention : on n'aura pas $t^p = s^p \circ \alpha^p$ ; il faudra donc
veiller à ne pas utiliser cette égalité ci-dessous.)

\medskip\noindent
Considérons
$$
J^\bullet = I^\bullet \times \prod\nolimits_p J^{p, \bullet}
$$
Comme les produits dans $D(\mathcal{A})$ s'obtiennent en prenant
les produits de complexes K-injectifs
(lemme \ref{lemma-derived-products}) et que $J^{p, \bullet}$
est isomorphe à $0$ dans $D(\mathcal{A})$, le morphisme
$J^\bullet \to I^\bullet$ est un isomorphisme dans $D(\mathcal{A})$.
Considérons l'application
$$
j = i \times (s^p \circ i)_{p \in \mathbf{Z}} :
K^\bullet
\longrightarrow
I^\bullet \times \prod\nolimits_p J^{p, \bullet} = J^\bullet
$$
D'après les remarques précédentes, c'est un quasi-isomorphisme. Il est aussi
injectif. Pour $p \in \mathbf{Z}$, on définit $F^pJ^\bullet \subset J^\bullet$ par
$$
\Im\left(
\alpha^p \times (t^{p'} \circ \alpha^{pp'})_{p' \leq p} :
I^{p, \bullet}
\to
I^\bullet \times \prod\nolimits_{p' \leq p} J^{p', \bullet}
\right)
\times \prod\nolimits_{p' > p} J^{p', \bullet}
$$
Ce complexe est isomorphe au complexe
$I^{p, \bullet} \times \prod_{p' > p} J^{p', \bullet}$
car $\alpha^{pp} = \text{id}$ et $t^p$ est injective.
Par conséquent, $F^pJ^\bullet$ est quasi-isomorphe à $I^{p, \bullet}$
(on raisonne comme ci-dessus). La commutativité du diagramme précédent donne
$j(F^pK^\bullet) \subset F^pJ^\bullet$. Le morphisme de complexes correspondant
$F^pK^\bullet \to F^pJ^\bullet$ est donc un quasi-isomorphisme.
Enfin, pour voir que
$F^{p + 1}J^\bullet \subset F^pJ^\bullet$
on utilise l'égalité $\alpha^{pp'} \circ \alpha^{p + 1p} = \alpha^{p + 1p'}$
et la commutativité du premier diagramme affiché au premier paragraphe
de la démonstration.

\medskip\noindent
Affirmons que $j : K^\bullet \to J^\bullet$ résout le problème posé par le
lemme. En effet, $F^pJ^n$ est un objet injectif de $\mathcal{A}$ puisqu'il est
isomorphe à $I^{p, n} \times \prod_{p' > p} J^{p', n}$ et que les produits
d'injectifs sont injectifs. L'application injective $F^pJ^n \to J^n$ est alors
scindée ; le quotient $J^n/F^pJ^n$ est donc lui aussi injectif, comme facteur
direct de l'objet injectif $J^n$. Il en va de même de
$F^pJ^n/F^{p'}J^n$. En particulier,
$0 \to F^pJ^\bullet \to J^\bullet \to J^\bullet/F^pJ^\bullet \to 0$
est une suite exacte courte de complexes scindée terme à terme ; elle définit
donc, par convention, un triangle distingué dans $K(\mathcal{A})$.
Puisque $J^\bullet$ et $F^pJ^\bullet$ sont des complexes K-injectifs, il en va
de même de $J^\bullet/F^pJ^\bullet$ d'après Catégories dérivées, lemme
\ref{derived-lemma-triangle-K-injective}. Un raisonnement semblable montre que
$F^pJ^\bullet/F^{p'}J^\bullet$ est K-injectif. Par construction,
$j : K^\bullet \to J^\bullet$ et les applications induites
$F^pK^\bullet \to F^pJ^\bullet$ sont des quasi-isomorphismes. Les longues
suites exactes de cohomologie des complexes considérés montrent qu'il en va
de même pour
$K^\bullet/F^pK^\bullet \to J^\bullet/F^pJ^\bullet$ et
$F^pK^\bullet/F^{p'}K^\bullet \to F^pJ^\bullet/F^{p'}J^\bullet$.
\end{proof}

\begin{remark}
\label{remark-ext-into-filtered-complex}
Soit $\mathcal{A}$ une catégorie abélienne de Grothendieck.
Soit $K^\bullet$ un complexe filtré de $\mathcal{A}$, voir
Homologie, définition \ref{homology-definition-filtered-complex}.
Pour alléger les notations, désignons par $K$, $F^pK$, $\text{gr}^pK$
les objets de $D(\mathcal{A})$ représentés respectivement par $K^\bullet$,
$F^pK^\bullet$, $\text{gr}^pK^\bullet$. Soit $M \in D(\mathcal{A})$.
Le lemme \ref{lemma-K-injective-embedding-filtration} permet de construire
une suite spectrale $(E_r, d_r)_{r \geq 1}$ d'objets bigradués de
$\mathcal{A}$, où $d_r$ est de bidegré $(r, -r + 1)$ et
$$
E_1^{p, q} = \Ext^{p + q}(M, \text{gr}^pK)
$$
Si, pour tout $n$, on a
$$
\Ext^n(M, F^pK) = 0 \text{ pour } p \gg 0
\quad\text{et}\quad
\Ext^n(M, F^pK) = \Ext^n(M, K) \text{ pour } p \ll 0
$$
alors la suite spectrale est bornée et converge vers $\Ext^{p + q}(M, K)$.
En effet, choisissons un complexe quelconque $M^\bullet$ représentant $M$,
un morphisme $j : K^\bullet \to J^\bullet$ comme dans le lemme, et considérons
le complexe
$$
\Hom^\bullet(M^\bullet, J^\bullet)
$$
défini exactement comme dans Compléments d'algèbre, section
\ref{more-algebra-section-hom-complexes}. En posant
$F^p\Hom^\bullet(M^\bullet, J^\bullet) =
\Hom^\bullet(M^\bullet, F^pJ^\bullet)$, on obtient un complexe filtré.
La suite spectrale d'Homologie, section
\ref{homology-section-filtered-complex}, possède les différentielles et les
termes décrits ci-dessus ; nous omettons les détails. Le caractère borné et la
convergence résultent d'Homologie, lemme
\ref{homology-lemma-ss-converges-trivial}.
\end{remark}

\begin{remark}
\label{remark-spectral-sequences-ext}
Soit $\mathcal{A}$ une catégorie abélienne de Grothendieck.
Soient $M, K$ des objets de $D(\mathcal{A})$.
Pour tout choix d'un complexe $K^\bullet$ représentant $K$, on peut utiliser
la filtration $F^pK^\bullet = \tau_{\leq -p}K^\bullet$ et la discussion de la
remarque \ref{remark-ext-into-filtered-complex} afin d'obtenir une suite
spectrale telle que
$$
E_1^{p, q} = \Ext^{2p + q}(M, H^{-p}(K))
$$
Cette suite spectrale ne dépend pas du choix du complexe $K^\bullet$
représentant $K$. Après la renumérotation $p = -j$ et $q = i + 2j$, on
obtient une suite spectrale $(E'_r, d'_r)_{r \geq 2}$ où $d'_r$ est de
bidegré $(r, -r + 1)$ et
$$
(E'_2)^{i, j} = \Ext^i(M, H^j(K))
$$
Si $M \in D^-(\mathcal{A})$ et $K \in D^+(\mathcal{A})$, alors $E_r$ et
$E'_r$ sont toutes deux bornées et convergent vers $\Ext^{p + q}(M, K)$.
Si l'on utilise la filtration $F^pK^\bullet = \sigma_{\geq p}K^\bullet$,
on obtient
$$
E_1^{p, q} = \Ext^q(M, K^p)
$$
Si $M \in D^-(\mathcal{A})$ et si $K^\bullet$ est borné inférieurement,
cette suite spectrale est bornée et converge vers $\Ext^{p + q}(M, K)$.
\end{remark}

\begin{remark}
\label{remark-ext-from-filtered-complex}
Soit $\mathcal{A}$ une catégorie abélienne de Grothendieck. Soit
$K \in D(\mathcal{A})$. Soit $M^\bullet$ un complexe filtré de
$\mathcal{A}$, voir Homologie, définition
\ref{homology-definition-filtered-complex}. Pour alléger les notations,
désignons par $M$, $M/F^pM$, $\text{gr}^pM$ les objets de
$D(\mathcal{A})$ représentés respectivement par $M^\bullet$,
$M^\bullet/F^pM^\bullet$, $\text{gr}^pM^\bullet$. Par dualité avec la
remarque \ref{remark-ext-into-filtered-complex}, on peut construire une suite
spectrale $(E_r, d_r)_{r \geq 1}$ d'objets bigradués de $\mathcal{A}$, où
$d_r$ est de bidegré $(r, -r + 1)$ et
$$
E_1^{p, q} = \Ext^{p + q}(\text{gr}^{-p}M, K)
$$
Si, pour tout $n$, on a
$$
\Ext^n(M/F^pM, K) = 0 \text{ pour } p \ll 0
\quad\text{et}\quad
\Ext^n(M/F^pM, K) = \Ext^n(M, K) \text{ pour } p \gg 0
$$
alors la suite spectrale est bornée et converge vers $\Ext^{p + q}(M, K)$.
En effet, choisissons un complexe K-injectif $I^\bullet$ à termes injectifs
représentant $K$, voir le théorème
\ref{theorem-K-injective-embedding-grothendieck}. Considérons le complexe
$$
\Hom^\bullet(M^\bullet, I^\bullet)
$$
défini exactement comme dans Compléments d'algèbre, section
\ref{more-algebra-section-hom-complexes}. En posant
$$
F^p\Hom^\bullet(M^\bullet, I^\bullet) =
\Hom^\bullet(M^\bullet/F^{-p + 1}M^\bullet, I^\bullet)
$$
on obtient un complexe filtré (remarquer le signe et le décalage de la
filtration). La suite spectrale d'Homologie, section
\ref{homology-section-filtered-complex}, possède les différentielles et les
termes décrits ci-dessus ; nous omettons les détails. Le caractère borné et la
convergence résultent d'Homologie, lemme
\ref{homology-lemma-ss-converges-trivial}.
\end{remark}

\begin{remark}
\label{remark-spectral-sequences-ext-variant}
Soit $\mathcal{A}$ une catégorie abélienne de Grothendieck.
Soient $M, K$ des objets de $D(\mathcal{A})$.
Pour tout choix d'un complexe $M^\bullet$ représentant $M$, on peut utiliser
la filtration $F^pM^\bullet = \tau_{\leq -p}M^\bullet$ et la discussion de la
remarque \ref{remark-ext-into-filtered-complex} afin d'obtenir une suite
spectrale telle que
$$
E_1^{p, q} = \Ext^{2p + q}(H^p(M), K)
$$
Cette suite spectrale ne dépend pas du choix du complexe $M^\bullet$
représentant $M$. Après la renumérotation $p = -j$ et $q = i + 2j$, on
obtient une suite spectrale $(E'_r, d'_r)_{r \geq 2}$ où $d'_r$ est de
bidegré $(r, -r + 1)$ et
$$
(E'_2)^{i, j} = \Ext^i(H^{-j}(M), K)
$$
Si $M \in D^-(\mathcal{A})$ et $K \in D^+(\mathcal{A})$, alors $E_r$ et
$E'_r$ sont bornées et convergent vers $\Ext^{p + q}(M, K)$.
Si l'on utilise la filtration $F^pM^\bullet = \sigma_{\geq p}M^\bullet$,
on obtient
$$
E_1^{p, q} = \Ext^q(M^{-p}, K)
$$
Si $K \in D^+(\mathcal{A})$ et si $M^\bullet$ est borné supérieurement,
cette suite spectrale est bornée et converge vers $\Ext^{p + q}(M, K)$.
\end{remark}

\begin{lemma}
\label{lemma-represent-by-filtered-complex}
Soit $\mathcal{A}$ une catégorie abélienne de Grothendieck. Supposons donnés
un objet $E \in D(\mathcal{A})$, un système projectif
$\{E^i\}_{i \in \mathbf{Z}}$ d'objets de $D(\mathcal{A})$ indexé par
$\mathbf{Z}$, ainsi qu'un système compatible de morphismes $E^i \to E$.
On a le diagramme suivant :
$$
\ldots \to E^{i + 1} \to E^i \to E^{i - 1} \to \ldots \to E
$$
Il existe alors un complexe filtré $K^\bullet$ de $\mathcal{A}$
(Homologie, définition \ref{homology-definition-filtered-complex}) tel que
$K^\bullet$ représente $E$ et que $F^iK^\bullet$ représente $E^i$, de façon
compatible avec les morphismes donnés.
\end{lemma}

\begin{proof}
D'après le théorème \ref{theorem-K-injective-embedding-grothendieck}, on peut
choisir un complexe K-injectif $I^\bullet$ représentant $E$, dont tous les
termes $I^n$ sont des objets injectifs de $\mathcal{A}$.
Choisissons un complexe $G^{0, \bullet}$ représentant $E^0$, puis un
morphisme de complexes $\varphi^0 : G^{0, \bullet} \to I^\bullet$
représentant $E^0 \to E$.
Pour $i > 0$, représentons par récurrence $E^i \to E^{i - 1}$ par un
morphisme de complexes
$\delta : G^{i, \bullet} \to G^{i - 1, \bullet}$
et posons $\varphi^i = \varphi^{i - 1} \circ \delta$.
Pour $i < 0$, représentons par récurrence $E^{i + 1} \to E^i$ par un
morphisme de complexes injectif terme à terme
$\delta : G^{i + 1, \bullet} \to G^{i, \bullet}$
(on peut, par exemple, utiliser Catégories dérivées, lemme
\ref{derived-lemma-make-injective}).
Affirmons qu'il existe un morphisme de complexes
$\varphi^i : G^{i, \bullet} \to I^\bullet$ représentant le morphisme
$E^i \to E$ et s'insérant dans le diagramme commutatif
$$
\xymatrix{
G^{i + 1, \bullet} \ar[r]_\delta \ar[d]_{\varphi^{i + 1}} &
G^{i, \bullet} \ar[ld]^{\varphi^i} \\
I^\bullet
}
$$
En effet, choisissons d'abord un morphisme de complexes quelconque
$\varphi : G^{i, \bullet} \to I^\bullet$ représentant le morphisme
$E^i \to E$. Les morphismes $\varphi \circ \delta$ et
$\varphi^{i + 1}$ sont alors homotopes par une homotopie
$h^p : G^{i + 1, p} \to I^{p - 1}$. Comme les termes de $I^\bullet$ sont
injectifs et que $\delta$ est injectif terme à terme, on peut prolonger
$h^p$ en $(h')^p : G^{i, p} \to I^{p - 1}$. En posant
$\varphi^i = \varphi + h' \circ d + d \circ h'$, on obtient ce qui était
annoncé.

\medskip\noindent
Choisissons ensuite, pour tout $i$, un morphisme de complexes injectif terme
à terme $a^i : G^{i, \bullet} \to J^{i, \bullet}$, où $J^{i, \bullet}$ est
acyclique et K-injectif et où les $J^{i, p}$ sont des objets injectifs de
$\mathcal{A}$. Pour cela, on envoie d'abord $G^{i, \bullet}$ dans le cône de
l'identité, puis on applique le théorème cité ci-dessus.
En raisonnant comme précédemment, on trouve des morphismes de complexes
$\delta' : J^{i, \bullet} \to J^{i - 1, \bullet}$ tels que les diagrammes
$$
\xymatrix{
G^{i, \bullet} \ar[r]_\delta \ar[d]_{a^i} &
G^{i - 1, \bullet} \ar[d]^{a^{i - 1}} \\
J^{i, \bullet} \ar[r]^{\delta'} &
J^{i - 1, \bullet}
}
$$
soient commutatifs. (On pourrait aussi utiliser la fonctorialité des cônes et
celle du théorème.) Considérons alors les morphismes
$$
\xymatrix{
G^{i + 1, \bullet} \times \prod\nolimits_{p > i + 1} J^{p, \bullet}
\ar[r] \ar[rd] &
G^{i, \bullet} \times \prod\nolimits_{p > i} J^{p, \bullet}
\ar[r] \ar[d] &
G^{i - 1, \bullet} \times \prod\nolimits_{p > i - 1} J^{p, \bullet}
\ar[ld] \\
& I^\bullet \times \prod\nolimits_p J^{p, \bullet}
}
$$
Ici, les flèches portant sur $J^{p, \bullet}$ sont les flèches évidentes
(l'identité ou le morphisme nul). Sur le facteur $G^{i, \bullet}$, on utilise
$\delta : G^{i, \bullet} \to G^{i - 1, \bullet}$, le morphisme
$\varphi^i : G^{i, \bullet} \to I^\bullet$, le morphisme nul
$0 : G^{i, \bullet} \to J^{p, \bullet}$ pour $p > i$, le morphisme
$a^i : G^{i, \bullet} \to J^{p, \bullet}$ pour $p = i$, et
$(\delta')^{i - p} \circ a^i = a^p \circ \delta^{i - p} :
G^{i, \bullet} \to J^{p, \bullet}$ pour $p < i$.
Nous omettons de vérifier que toutes les flèches du diagramme sont injectives
terme à terme. On obtient ainsi un complexe filtré. Comme les produits dans
$D(\mathcal{A})$ s'obtiennent en prenant les produits de complexes K-injectifs
(lemme \ref{lemma-derived-products}) et que $J^{p, \bullet}$ est nul dans
$D(\mathcal{A})$, ce diagramme représente le diagramme donné dans la catégorie
dérivée. Cela achève la démonstration.
\end{proof}

\begin{lemma}
\label{lemma-represent-by-filtered-complex-bis}
Dans la situation du lemme \ref{lemma-represent-by-filtered-complex},
supposons donné un second système projectif
$\{(E')^i\}_{i \in \mathbf{Z}}$, ainsi qu'un système compatible de morphismes
$(E')^i \to E$. Il existe alors un complexe bifiltré $K^\bullet$ de
$\mathcal{A}$ tel que $K^\bullet$ représente $E$, que $F^iK^\bullet$
représente $E^i$ et que $(F')^iK^\bullet$ représente $(E')^i$, de façon
compatible avec les morphismes donnés.
\end{lemma}

\begin{proof}
Le lemme permet d'abord de choisir $K^\bullet$ et $F$. On choisit ensuite
$(K')^\bullet$ et $F'$ convenant à $\{(E')^i\}_{i \in \mathbf{Z}}$ et aux
morphismes $(E')^i \to E$. D'après le lemme
\ref{lemma-K-injective-embedding-filtration}, on peut supposer que
$K^\bullet$ est un complexe K-injectif. On choisit alors un morphisme de
complexes $(K')^\bullet \to K^\bullet$ correspondant aux identifications
données $(K')^\bullet \cong E \cong K^\bullet$. On peut en outre choisir un
morphisme $(K')^\bullet \to J^\bullet$ injectif terme à terme, où
$J^\bullet$ est acyclique et K-injectif. (Pour cela, on envoie d'abord
$(K')^\bullet$ dans le cône de l'identité, puis on applique le théorème
\ref{theorem-K-injective-embedding-grothendieck}.) Les morphismes
$(K')^\bullet \to K^\bullet \times J^\bullet$ et
$K^\bullet \to K^\bullet \times J^\bullet$ sont alors tous deux injectifs
terme à terme et sont des quasi-isomorphismes (car le produit représente $E$
d'après le lemme \ref{lemma-derived-products}). Il suffit enfin de prendre les
images des filtrations de $K^\bullet$ et $(K')^\bullet$ par ces morphismes.
\end{proof}





\section{Le théorème de Gabriel--Popescu}
\label{section-gabriel-popescu}

\noindent
Dans cette section, nous étudions le théorème principal de \cite{GP}. La
méthode de démonstration suit des notes de Jacob Lurie et d'Akhil Mathew, qui
suivent à leur tour la présentation de Kuhn dans \cite{Kuhn}. Voir aussi
\cite{Takeuchi}.

\medskip\noindent
Soit $\mathcal{A}$ une catégorie abélienne de Grothendieck et soit $U$ un
générateur de $\mathcal{A}$, voir la définition
\ref{definition-grothendieck-conditions}. Posons
$R = \Hom_\mathcal{A}(U, U)$. Considérons le foncteur
$G : \mathcal{A} \to \text{Mod}_R$ défini par
$$
G(A) = \Hom_\mathcal{A}(U, A)
$$
muni de sa structure canonique de $R$-module à droite.

\begin{lemma}
\label{lemma-gabriel-popescu-left-adjoint}
Le foncteur $G$ ci-dessus possède un adjoint à gauche
$F : \text{Mod}_R \to \mathcal{A}$.
\end{lemma}

\begin{proof}
Nous donnons deux démonstrations de ce lemme.

\medskip\noindent
La première utilise le théorème du foncteur adjoint, voir Catégories,
théorème \ref{categories-theorem-adjoint-functor}. Remarquons que
$G : \mathcal{A} \to \text{Mod}_R$ est exact à gauche et envoie les produits
sur les produits. Le foncteur $G$ commute donc aux limites. Pour vérifier la condition
ensembliste du théorème, soit $M$ un objet de $\text{Mod}_R$. Choisissons un
cardinal $\kappa$ suffisamment grand et désignons par $E$ un ensemble d'objets
de $\mathcal{A}$ tel que tout objet $A$ vérifiant $|A| \leq \kappa$ soit
isomorphe à un élément de $E$. C'est possible d'après le lemme
\ref{lemma-set-iso-classes-bounded-size}. Posons
$I = \coprod_{A \in E} \Hom_R(M, G(A))$. Un élément $i \in I$ sera considéré
comme un couple $(A_i, f_i)$. Enfin, soient $A$ un objet quelconque de
$\mathcal{A}$ et $f : M \to G(A)$ un morphisme quelconque. Interprétons les
éléments de $\Im(f) \subset G(A) = \Hom_\mathcal{A}(U, A)$ comme des
morphismes $u : U \to A$. Posons
$$
A' = \Im(\bigoplus\nolimits_{u \in \Im(f)} U \xrightarrow{u} A)
$$
Puisque $G$ est exact à gauche, $G(A') \subset G(A)$ contient $\Im(f)$ ; on
obtient donc $f' : M \to G(A')$ par lequel $f$ se factorise. D'autre part,
l'objet $A'$ est quotient d'une somme directe d'au plus $|M|$ copies de $U$.
Ainsi, si $\kappa = |\bigoplus_{|M|} U|$, le couple $(A', f')$ est isomorphe
à un élément $(A_i, f_i)$ de $I$, et $f$ se factorise, comme souhaité, sous
la forme $M \xrightarrow{f_i} G(A_i) \to G(A)$.

\medskip\noindent
La seconde démonstration construit $F$ et montre, en un certain sens, que
``$F(M) = M \otimes_R U$''. En effet, pour tout $R$-module $M$, on peut
choisir une présentation
$$
\bigoplus\nolimits_{j \in J} R \to
\bigoplus\nolimits_{i \in I} R \to
M \to 0
$$
On définit alors $F(M)$ par la suite exacte correspondante
$$
\bigoplus\nolimits_{j \in J} U \to
\bigoplus\nolimits_{i \in I} U \to
F(M) \to 0
$$
Cette construction ne dépend pas du choix de la présentation et elle est
fonctorielle ; nous omettons les détails. Pour tout $A$ dans $\mathcal{A}$,
on obtient une suite exacte
$$
0 \to \Hom_\mathcal{A}(F(M), A) \to
\prod\nolimits_{i \in I} G(A) \to
\prod\nolimits_{j \in J} G(A)
$$
qui est isomorphe à la suite
$$
0 \to \Hom_R(M, G(A)) \to
\Hom_R(\bigoplus\nolimits_{i \in I} R, G(A)) \to
\Hom_R(\bigoplus\nolimits_{j \in J} R, G(A))
$$
ce qui montre que $F$ est l'adjoint à gauche de $G$.
\end{proof}

\begin{lemma}
\label{lemma-F-G-monos}
Soit $f : M \to G(A)$ un morphisme injectif dans $\text{Mod}_R$.
Alors le morphisme adjoint $f' : F(M) \to A$ est lui aussi injectif.
\end{lemma}

\begin{proof}
Choisissons un morphisme $R^{\oplus n} \to M$ et considérons le morphisme
correspondant $U^{\oplus n} \to F(M)$. Soit
$v : U \to U^{\oplus n}$ un morphisme tel que la composée
$U \to U^{\oplus n} \to F(M) \to A$ soit $0$. La flèche
$v : U \to U^{\oplus n}$ est alors un élément $v$ de $R^{\oplus n}$ dont
l'image dans $G(A)$ est nulle. Puisque $f$ est injectif, l'image de $v$ dans
$M$ est nulle ; cela signifie que $U \to U^{\oplus n} \to F(M)$ est nul,
par la construction de $F(M)$ donnée dans la démonstration du
lemme \ref{lemma-gabriel-popescu-left-adjoint}. Comme $U$ est un générateur, on en
déduit que
$$
\Ker(U^{\oplus n} \to F(M) \to A) = \Ker(U^{\oplus n} \to F(M))
$$
Pour achever la démonstration, choisissons une surjection
$\bigoplus_{i \in I} R \to M$ et considérons la surjection correspondante
$$
\pi : \bigoplus\nolimits_{i \in I} U \longrightarrow F(M)
$$
Pour montrer que $f'$ est injectif, il suffit d'établir l'égalité
$\Ker(\pi) = \Ker(f' \circ \pi)$ entre sous-objets de
$\bigoplus_{i \in I} U$. Or on peut écrire $\bigoplus_{i \in I} U$ comme la
colimite filtrante de ses sous-objets $\bigoplus_{i \in I'} U$, où
$I' \subset I$ parcourt les sous-ensembles finis. Comme les colimites
filtrantes sont exactes dans $\mathcal{A}$ par AB5, on a
$$
\Ker(\pi) =
\colim_{I' \subset I\text{ fini}}
\left(\bigoplus\nolimits_{i \in I'} U\right)
\bigcap \Ker(\pi)
$$
et
$$
\Ker(f' \circ \pi) =
\colim_{I' \subset I\text{ fini}}
\left(\bigoplus\nolimits_{i \in I'} U\right)
\bigcap \Ker(f' \circ \pi)
$$
et l'on obtient l'égalité, car la première égalité affichée ci-dessus
l'établit pour chaque $I'$.
\end{proof}

\begin{theorem}
\label{theorem-gabriel-popescu}
Soit $\mathcal{A}$ une catégorie abélienne de Grothendieck. Il existe alors
un anneau (non commutatif) $R$ et des foncteurs
$G : \mathcal{A} \to \text{Mod}_R$ et
$F : \text{Mod}_R \to \mathcal{A}$ tels que
\begin{enumerate}
\item $F$ est l'adjoint à gauche de $G$,
\item $G$ est pleinement fidèle, et
\item $F$ est exact.
\end{enumerate}
Ces foncteurs sont en outre ceux construits ci-dessus.
\end{theorem}

\begin{proof}
Montrons d'abord que $G$ est pleinement fidèle, ou, ce qui revient au même,
que $F \circ G \to \text{id}$ est un isomorphisme, voir Catégories, lemme
\ref{categories-lemma-adjoint-fully-faithful}. Pour tout objet $A$, le
morphisme $F(G(A)) \to A$ est surjectif, car, par construction, tout morphisme
$U \to A$ se factorise par $F(G(A))$. D'autre part, le morphisme
$F(G(A)) \to A$ est l'adjoint du morphisme
$\text{id} : G(A) \to G(A)$ ; il est donc injectif d'après le lemme
\ref{lemma-F-G-monos}.

\medskip\noindent
Le foncteur $F$ est exact à droite puisqu'il est adjoint à gauche.
Comme $\text{Mod}_R$ possède assez de projectifs, pour montrer que $F$ est
exact, il suffit de montrer que son premier foncteur dérivé gauche $L_1F$ est
nul. Pour établir $L_1F(M) = 0$ pour un $R$-module $M$, choisissons une suite
exacte $0 \to K \to P \to M \to 0$ de $R$-modules où $P$ est libre. Il suffit
de montrer que $F(K) \to F(P)$ est injectif. Écrivons cette suite comme
colimite filtrante de suites $0 \to K_i \to P_i \to M_i \to 0$ où $P_i$ est
un $R$-module libre de type fini : il suffit d'écrire ainsi $P$, puis de poser
$K_i = K \cap P_i$ et $M_i = \Im(P_i \to M)$. Puisque $F$ est adjoint à
gauche, il commute aux colimites ; et puisque $\mathcal{A}$ est une catégorie
abélienne de Grothendieck, $F(K) \to F(P)$ est injectif dès que chaque
$F(K_i) \to F(P_i)$ l'est. Il suffit donc de vérifier que
$F(K) \to F(P)$ est injectif lorsque $K \subset P = R^{\oplus n}$.
Le lemme \ref{lemma-F-G-monos} montre précisément que
$F(K) \to U^{\oplus n}$ est injectif.
\end{proof}

\begin{lemma}
\label{lemma-gabriel-popescu}
\begin{reference}
\cite[Corollary 4.1]{serpe}
\end{reference}
Soit $\mathcal{A}$ une catégorie abélienne de Grothendieck. Soient
$R$, $F$, $G$ comme dans le théorème de Gabriel--Popescu
(théorème \ref{theorem-gabriel-popescu}). On obtient alors les foncteurs
dérivés
$$
RG : D(\mathcal{A}) \to D(\text{Mod}_R)
\quad\text{et}\quad
F : D(\text{Mod}_R) \to D(\mathcal{A})
$$
tels que $F$ est adjoint à gauche de $RG$, que $RG$ est pleinement fidèle et
que $F \circ RG = \text{id}$.
\end{lemma}

\begin{proof}
L'existence et l'adjonction de ces foncteurs résultent des théorèmes
\ref{theorem-gabriel-popescu} et
\ref{theorem-K-injective-embedding-grothendieck}
et
des lemmes de Catégories dérivées
\ref{derived-lemma-K-injective-defined},
\ref{derived-lemma-right-derived-exact-functor}, et
\ref{derived-lemma-derived-adjoint-functors}.
L'égalité $F \circ RG = \text{id}$ vient du fait qu'on peut calculer $RG$ sur
un objet de $D(\mathcal{A})$ en appliquant $G$ à un complexe représentant
convenable $I^\bullet$ (par exemple un complexe K-injectif) ; on a alors
$F(G(I^\bullet)) = I^\bullet$ puisque $F \circ G = \text{id}$. La pleine
fidélité de $RG$ en résulte d'après Catégories, lemme
\ref{categories-lemma-adjoint-fully-faithful}.
\end{proof}





\section{Représentabilité de Brown et catégories abéliennes de Grothendieck}
\label{section-brown}

\noindent
Dans cette section, nous démontrons rapidement un théorème de représentabilité
pour les catégories dérivées de catégories abéliennes de Grothendieck. Le
lecteur devrait d'abord étudier le cas des catégories triangulées compactement
engendrées dans Catégories dérivées, section \ref{derived-section-brown}.
Il est ensuite tout indiqué de consulter, plutôt que cette section, la
littérature pour des résultats plus généraux de ce type ; voir par exemple
\cite{Franke}, \cite{Neeman}, \cite{Krause}, ou Catégories dérivées, section
\ref{derived-section-brown-bis}.

\begin{lemma}
\label{lemma-brown}
Soit $\mathcal{A}$ une catégorie abélienne de Grothendieck.
Soit $H : D(\mathcal{A}) \to \textit{Ab}$ un foncteur cohomologique
contravariant qui transforme les sommes directes en produits.
Alors $H$ est représentable.
\end{lemma}

\begin{proof}
Soient $R, F, G, RG$ comme dans le lemme \ref{lemma-gabriel-popescu}, et
considérons le foncteur $H \circ F : D(\text{Mod}_R) \to \textit{Ab}$.
Comme $F$ est adjoint à gauche, il envoie les sommes directes sur les sommes
directes ; $H \circ F$ transforme donc les sommes directes en produits.
D'autre part, la catégorie dérivée $D(\text{Mod}_R)$ est engendrée par un
unique objet compact, à savoir $R$. D'après Catégories dérivées, lemme
\ref{derived-lemma-brown}, le foncteur $H \circ F$ est représentable, disons
par $L \in D(\text{Mod}_R)$. Choisissons un triangle distingué
$$
M \to L \to RG(F(L)) \to M[1]
$$
dans $D(\text{Mod}_R)$. Alors $F(M) = 0$ puisque
$F \circ RG = \text{id}$. Par conséquent $H(F(M)) = 0$, puis
$\Hom(M, L) = 0$. Il s'ensuit que $L \to RG(F(L))$ est l'inclusion d'un
facteur direct, voir Catégories dérivées, lemme \ref{derived-lemma-split}.
Pour $A$ dans $D(\mathcal{A})$, on obtient
\begin{align*}
H(A)
& =
H(F(RG(A))) \\
& =
\Hom(RG(A), L) \\
& \to
\Hom(RG(A), RG(F(L))) \\
& =
\Hom(F(RG(A)), F(L)) \\
& =
\Hom(A, F(L))
\end{align*}
où la flèche possède un inverse à gauche fonctoriel en $A$. Autrement dit,
$H$ est facteur direct d'un foncteur représentable. Comme $D(\mathcal{A})$
est karoubienne (Catégories dérivées, lemme
\ref{derived-lemma-projectors-have-images-triangulated}), on conclut.
\end{proof}

\begin{proposition}
\label{proposition-brown}
Soit $\mathcal{A}$ une catégorie abélienne de Grothendieck. Soit $\mathcal{D}$
une catégorie triangulée. Soit $F : D(\mathcal{A}) \to \mathcal{D}$ un
foncteur exact de catégories triangulées qui transforme les sommes directes en
sommes directes. Alors $F$ possède un adjoint à droite exact.
\end{proposition}

\begin{proof}
Pour un objet $Y$ de $\mathcal{D}$, considérons le foncteur contravariant
$$
D(\mathcal{A}) \to \textit{Ab},\quad
W \mapsto \Hom_\mathcal{D}(F(W), Y)
$$
C'est un foncteur cohomologique puisque $F$ est exact, et il transforme les
sommes directes en produits puisque $F$ transforme les sommes directes en
sommes directes. Le lemme \ref{lemma-brown} fournit donc un objet $X$ de
$D(\mathcal{A})$ tel que
$\Hom_{D(\mathcal{A})}(W, X) = \Hom_\mathcal{D}(F(W), Y)$.
L'existence de l'adjoint résulte de Catégories, lemme
\ref{categories-lemma-adjoint-exists}. Son exactitude résulte de Catégories
dérivées, lemme \ref{derived-lemma-adjoint-is-exact}.
\end{proof}











\begin{multicols}{2}[\section{Autres chapitres}]
\noindent
Préliminaires
\begin{enumerate}
\item \hyperref[introduction-section-phantom]{Introduction}
\item \hyperref[conventions-section-phantom]{Conventions}
\item \hyperref[sets-section-phantom]{Théorie des ensembles}
\item \hyperref[categories-section-phantom]{Catégories}
\item \hyperref[topology-section-phantom]{Topologie}
\item \hyperref[sheaves-section-phantom]{Faisceaux sur les espaces}
\item \hyperref[sites-section-phantom]{Sites et faisceaux}
\item \hyperref[stacks-section-phantom]{Champs}
\item \hyperref[fields-section-phantom]{Corps}
\item \hyperref[algebra-section-phantom]{Algèbre commutative}
\item \hyperref[brauer-section-phantom]{Groupes de Brauer}
\item \hyperref[homology-section-phantom]{Algèbre homologique}
\item \hyperref[derived-section-phantom]{Catégories dérivées}
\item \hyperref[simplicial-section-phantom]{Méthodes simpliciales}
\item \hyperref[more-algebra-section-phantom]{Compléments d'algèbre}
\item \hyperref[smoothing-section-phantom]{Lissage des morphismes d'anneaux}
\item \hyperref[modules-section-phantom]{Faisceaux de modules}
\item \hyperref[sites-modules-section-phantom]{Modules sur les sites}
\item \hyperref[injectives-section-phantom]{Objets injectifs}
\item \hyperref[cohomology-section-phantom]{Cohomologie des faisceaux}
\item \hyperref[sites-cohomology-section-phantom]{Cohomologie sur les sites}
\item \hyperref[dga-section-phantom]{Algèbre différentielle graduée}
\item \hyperref[dpa-section-phantom]{Algèbre à puissances divisées}
\item \hyperref[sdga-section-phantom]{Faisceaux différentiels gradués}
\item \hyperref[hypercovering-section-phantom]{Hyperrecouvrements}
\end{enumerate}
Schémas
\begin{enumerate}
\setcounter{enumi}{25}
\item \hyperref[schemes-section-phantom]{Schémas}
\item \hyperref[constructions-section-phantom]{Constructions de schémas}
\item \hyperref[properties-section-phantom]{Propriétés des schémas}
\item \hyperref[morphisms-section-phantom]{Morphismes de schémas}
\item \hyperref[coherent-section-phantom]{Cohomologie des schémas}
\item \hyperref[divisors-section-phantom]{Diviseurs}
\item \hyperref[limits-section-phantom]{Limites de schémas}
\item \hyperref[varieties-section-phantom]{Variétés}
\item \hyperref[topologies-section-phantom]{Topologies sur les schémas}
\item \hyperref[descent-section-phantom]{Descente}
\item \hyperref[perfect-section-phantom]{Catégories dérivées de schémas}
\item \hyperref[more-morphisms-section-phantom]{Compléments sur les morphismes}
\item \hyperref[flat-section-phantom]{Compléments sur la platitude}
\item \hyperref[groupoids-section-phantom]{Schémas en groupoïdes}
\item \hyperref[more-groupoids-section-phantom]{Compléments sur les schémas en groupoïdes}
\item \hyperref[etale-section-phantom]{Morphismes étales de schémas}
\end{enumerate}
Thèmes de théorie des schémas
\begin{enumerate}
\setcounter{enumi}{41}
\item \hyperref[chow-section-phantom]{Homologie de Chow}
\item \hyperref[intersection-section-phantom]{Théorie de l'intersection}
\item \hyperref[pic-section-phantom]{Schémas de Picard des courbes}
\item \hyperref[weil-section-phantom]{Théories cohomologiques de Weil}
\item \hyperref[adequate-section-phantom]{Modules adéquats}
\item \hyperref[dualizing-section-phantom]{Complexes dualisants}
\item \hyperref[duality-section-phantom]{Dualité pour les schémas}
\item \hyperref[discriminant-section-phantom]{Discriminants et différentes}
\item \hyperref[derham-section-phantom]{Cohomologie de de Rham}
\item \hyperref[local-cohomology-section-phantom]{Cohomologie locale}
\item \hyperref[algebraization-section-phantom]{Géométrie algébrique et formelle}
\item \hyperref[curves-section-phantom]{Courbes algébriques}
\item \hyperref[resolve-section-phantom]{Résolution des surfaces}
\item \hyperref[models-section-phantom]{Réduction semi-stable}
\item \hyperref[functors-section-phantom]{Foncteurs et morphismes}
\item \hyperref[equiv-section-phantom]{Catégories dérivées de variétés}
\item \hyperref[pione-section-phantom]{Groupes fondamentaux de schémas}
\item \hyperref[etale-cohomology-section-phantom]{Cohomologie étale}
\item \hyperref[crystalline-section-phantom]{Cohomologie cristalline}
\item \hyperref[proetale-section-phantom]{Cohomologie pro-étale}
\item \hyperref[relative-cycles-section-phantom]{Cycles relatifs}
\item \hyperref[more-etale-section-phantom]{Compléments de cohomologie étale}
\item \hyperref[trace-section-phantom]{La formule des traces}
\end{enumerate}
Espaces algébriques
\begin{enumerate}
\setcounter{enumi}{64}
\item \hyperref[spaces-section-phantom]{Espaces algébriques}
\item \hyperref[spaces-properties-section-phantom]{Propriétés des espaces algébriques}
\item \hyperref[spaces-morphisms-section-phantom]{Morphismes d'espaces algébriques}
\item \hyperref[decent-spaces-section-phantom]{Espaces algébriques décents}
\item \hyperref[spaces-cohomology-section-phantom]{Cohomologie des espaces algébriques}
\item \hyperref[spaces-limits-section-phantom]{Limites d'espaces algébriques}
\item \hyperref[spaces-divisors-section-phantom]{Diviseurs sur les espaces algébriques}
\item \hyperref[spaces-over-fields-section-phantom]{Espaces algébriques sur des corps}
\item \hyperref[spaces-topologies-section-phantom]{Topologies sur les espaces algébriques}
\item \hyperref[spaces-descent-section-phantom]{Descente et espaces algébriques}
\item \hyperref[spaces-perfect-section-phantom]{Catégories dérivées d'espaces}
\item \hyperref[spaces-more-morphisms-section-phantom]{Compléments sur les morphismes d'espaces}
\item \hyperref[spaces-flat-section-phantom]{Platitude sur les espaces algébriques}
\item \hyperref[spaces-groupoids-section-phantom]{Groupoïdes dans les espaces algébriques}
\item \hyperref[spaces-more-groupoids-section-phantom]{Compléments sur les groupoïdes dans les espaces}
\item \hyperref[bootstrap-section-phantom]{Amorçage}
\item \hyperref[spaces-pushouts-section-phantom]{Sommes amalgamées d'espaces algébriques}
\end{enumerate}
Thèmes de géométrie
\begin{enumerate}
\setcounter{enumi}{81}
\item \hyperref[spaces-chow-section-phantom]{Groupes de Chow des espaces}
\item \hyperref[groupoids-quotients-section-phantom]{Quotients de groupoïdes}
\item \hyperref[spaces-more-cohomology-section-phantom]{Compléments sur la cohomologie des espaces}
\item \hyperref[spaces-simplicial-section-phantom]{Espaces simpliciaux}
\item \hyperref[spaces-duality-section-phantom]{Dualité pour les espaces}
\item \hyperref[formal-spaces-section-phantom]{Espaces algébriques formels}
\item \hyperref[restricted-section-phantom]{Algébrisation des espaces formels}
\item \hyperref[spaces-resolve-section-phantom]{Résolution des surfaces revisitée}
\end{enumerate}
Théorie des déformations
\begin{enumerate}
\setcounter{enumi}{89}
\item \hyperref[formal-defos-section-phantom]{Théorie formelle des déformations}
\item \hyperref[defos-section-phantom]{Théorie des déformations}
\item \hyperref[cotangent-section-phantom]{Le complexe cotangent}
\item \hyperref[examples-defos-section-phantom]{Problèmes de déformation}
\end{enumerate}
Champs algébriques
\begin{enumerate}
\setcounter{enumi}{93}
\item \hyperref[algebraic-section-phantom]{Champs algébriques}
\item \hyperref[examples-stacks-section-phantom]{Exemples de champs}
\item \hyperref[stacks-sheaves-section-phantom]{Faisceaux sur les champs algébriques}
\item \hyperref[criteria-section-phantom]{Critères de représentabilité}
\item \hyperref[artin-section-phantom]{Axiomes d'Artin}
\item \hyperref[quot-section-phantom]{Espaces de Quot et de Hilbert}
\item \hyperref[stacks-properties-section-phantom]{Propriétés des champs algébriques}
\item \hyperref[stacks-morphisms-section-phantom]{Morphismes de champs algébriques}
\item \hyperref[stacks-limits-section-phantom]{Limites de champs algébriques}
\item \hyperref[stacks-cohomology-section-phantom]{Cohomologie des champs algébriques}
\item \hyperref[stacks-perfect-section-phantom]{Catégories dérivées de champs}
\item \hyperref[stacks-introduction-section-phantom]{Introduction aux champs algébriques}
\item \hyperref[stacks-more-morphisms-section-phantom]{Compléments sur les morphismes de champs}
\item \hyperref[stacks-geometry-section-phantom]{La géométrie des champs}
\end{enumerate}
Thèmes de théorie des espaces de modules
\begin{enumerate}
\setcounter{enumi}{107}
\item \hyperref[moduli-section-phantom]{Champs de modules}
\item \hyperref[moduli-curves-section-phantom]{Espaces de modules de courbes}
\end{enumerate}
Divers
\begin{enumerate}
\setcounter{enumi}{109}
\item \hyperref[examples-section-phantom]{Exemples}
\item \hyperref[exercises-section-phantom]{Exercices}
\item \hyperref[guide-section-phantom]{Guide bibliographique}
\item \hyperref[desirables-section-phantom]{Desiderata}
\item \hyperref[coding-section-phantom]{Style de programmation}
\item \hyperref[obsolete-section-phantom]{Obsolète}
\item \hyperref[fdl-section-phantom]{Licence de documentation libre GNU}
\item \hyperref[index-section-phantom]{Index généré automatiquement}
\end{enumerate}
\end{multicols}



\bibliography{my}
\bibliographystyle{amsalpha}

\end{document}
